\documentclass[11pt,a4paper]{article}
\usepackage[T1]{fontenc}
\usepackage[utf8]{inputenc}
\usepackage[polish]{babel}
\usepackage[margin=2.3cm]{geometry}
\usepackage{lmodern,microtype}
\usepackage{array,longtable,booktabs,tabularx}
\usepackage{amsmath,amssymb}
\usepackage{graphicx,xcolor,tikz,listings}
\usetikzlibrary{arrows.meta,positioning,fit,backgrounds,shapes.geometric,shapes.multipart}
\usepackage[hidelinks,unicode]{hyperref}
\usepackage{xurl}
\usepackage{bookmark}
\usepackage{enumitem}

\definecolor{accent}{HTML}{245777}
\definecolor{pale}{HTML}{EDF3F7}
\definecolor{codebg}{HTML}{F5F5F5}
\hypersetup{
  pdftitle={TSDB -- architektura i szczegółowy projekt techniczny},
  pdfauthor={Hubert Błonowski, Jan Zaborski}
}
\setlength{\emergencystretch}{2em}
\setlength{\parskip}{0.35em}
\setlength{\parindent}{0pt}
\setcounter{tocdepth}{1}
\setlist{itemsep=0.15em,topsep=0.3em}
\renewcommand{\arraystretch}{1.18}
\setlength{\tabcolsep}{2.8pt}
\newcolumntype{P}[1]{>{\raggedright\arraybackslash}p{#1}}
\newcolumntype{Y}{>{\raggedright\arraybackslash}X}
\newcommand{\code}[1]{\expandafter\nolinkurl\expandafter{\detokenize{#1}}}
\newcommand{\tabhead}[1]{\toprule #1\\ \midrule\endfirsthead
  \toprule #1\\ \midrule\endhead \bottomrule\endfoot}
\lstset{
  basicstyle=\ttfamily\footnotesize,
  backgroundcolor=\color{codebg},
  frame=single,rulecolor=\color{black!15},
  breaklines=true,columns=fullflexible,keepspaces=true,
  showstringspaces=false,tabsize=2,
  aboveskip=0.8em,belowskip=0.8em
}
\tikzset{
  box/.style={draw=accent,rounded corners=2pt,fill=pale,
    align=center,font=\small,minimum height=11mm,text width=36mm},
  external/.style={box,fill=white,draw=black!60},
  arrow/.style={-{Latex[length=2mm]},thick,draw=accent},
  note/.style={font=\scriptsize,align=center,fill=white,inner sep=2pt},
  state/.style={box,text width=28mm,minimum height=10mm},
  umlclass/.style={draw=accent,fill=pale,rectangle split,
    rectangle split parts=3,rectangle split part align={left,left,left},
    align=left,font=\scriptsize,
    inner sep=5pt,text width=43mm},
  assoc/.style={thick,draw=accent},
  compose/.style={{Diamond[fill=accent,length=2mm]}-,thick,draw=accent},
  usecase/.style={ellipse,draw=accent,fill=pale,align=center,
    font=\small,text width=36mm,minimum height=13mm}
}

\begin{document}
\raggedbottom
\hypersetup{pageanchor=false}
\begin{titlepage}
\centering
{\large Politechnika Warszawska\\
Wydział Matematyki i Nauk Informacyjnych\par}
\vspace{2.1cm}
{\Large Projekt zespołowy -- laboratorium 2\par}
\vspace{1.2cm}
{\huge\bfseries Implementacja bazy danych\\
szeregów czasowych\\w języku Rust\par}
\vspace{1.1cm}
{\LARGE Architektura systemu\\i szczegółowy projekt techniczny\par}
\vspace{1.4cm}
\begin{tabular}{rl}
Autorzy: & Hubert Błonowski\\
         & Jan Zaborski\\[0.5em]
Promotor: & mgr inż. Paweł Sobótka\\[0.5em]
Wersja dokumentu: & 1.6\\
Data: & 3 listopada 2026 r.
\end{tabular}
\vfill
\url{https://github.com/janzaborski/rust-tsdb}
\vspace{0.8cm}
\end{titlepage}
\hypersetup{pageanchor=true}
\pagenumbering{roman}

\section*{Streszczenie}
\addcontentsline{toc}{section}{Streszczenie}
Projekt jednowęzłowej TSDB w Rust obejmuje moduły, model danych,
HTTP, podzbiór PromQL i Grafanę. Jeden Index z osobnymi blokadami Series
pozwala współbieżnie zapisywać, a rotacja bloków zachowuje dostępność
podczas flush. WAL z gałęzi \code{23-wal-improvements} zachowuje bufor RAM,
flusher, okresowy syncer oraz tryby None i Written. Checkpoint
wykorzystuje cut() i granicę całych plików. Kontrakty, formaty, indeksy i recovery umożliwiają
niezależną implementację magazynu, zapytań i HTTP.
To \textbf{specyfikacja docelowa}, nie opis prototypu ani wyniki pomiarów.
Punktem wyjścia jest lab1; doprecyzowania wymagań wskazano jawnie.

\section*{Historia zmian}
\addcontentsline{toc}{section}{Historia zmian}
\begin{tabularx}{\linewidth}{P{14mm}P{22mm}P{42mm}Y}
\toprule
Wersja & Data & Autorzy & Zakres zmiany\\ \midrule
1.0 & 24.10.2026 & H. Błonowski,\newline J. Zaborski &
Architektura, API, formaty, diagramy, GUI i plan weryfikacji.\\
1.1 & 24.10.2026 & H. Błonowski,\newline J. Zaborski &
Index/Series, współbieżność bez shardów, rotacja bloków i flush w tle.\\
1.2 & 24.10.2026 & H. Błonowski,\newline J. Zaborski &
Buforowany WAL, bajtowy LSN, poziomy trwałości, checkpoint i recovery.\\
1.3 & 02.11.2026 & H. Błonowski,\newline J. Zaborski &
Krótszy opis; UML przypadków użycia, klas, wdrożenia i stanów.\\
1.4 & 02.11.2026 & H. Błonowski,\newline J. Zaborski &
WAL z brancha 23-wal-improvements; dostosowanie head, checkpointu,
recovery, formatów i API.\\
1.6 & 03.11.2026 & H. Błonowski,\newline J. Zaborski &
Ujednolicenie nazw typów, pól, metod i tras HTTP z kodem.\\
\bottomrule
\end{tabularx}
\par
\begingroup
\footnotesize
\setlength{\parskip}{0pt}
\makeatletter
\renewcommand*\l@section[2]{%
  \@dottedtocline{1}{0em}{1.8em}{\bfseries #1}{#2}}
\makeatother
\tableofcontents
\endgroup
\clearpage
\pagenumbering{arabic}

\section{Cel, zakres i decyzje architektoniczne}
\label{sec:scope}

\subsection{Cel, odbiorcy i granice systemu}
Baza przyjmuje pomiary aplikacji, przechowuje je we własnym silniku
i udostępnia historię oraz agregacje w Grafanie. Projekt dla autorów,
promotora i recenzenta określa zachowanie docelowe; przykłady i cele
wydajnościowe nie są wynikami prototypu.

Zakres: jednowęzłowy silnik Rust, HTTP, dokładne dopasowanie etykiet,
pięć agregacji czasowych i demonstracja w Compose. Grafana i generatory
są klientami zewnętrznymi. Poza zakresem pozostają klaster, replikacja,
kompresja, retencja, kompaktowanie, alertowanie, pobieranie
\code{/metrics} i własny frontend. Dane rosną do usunięcia środowiska;
kopia zapasowa wymaga zatrzymania bazy i skopiowania całego katalogu.

\subsection{Doprecyzowania względem laboratorium 1}
\begin{longtable}{P{23mm}P{133mm}}
\tabhead{Wymaganie & Przyjęta interpretacja i konsekwencja}
FR-01 & Atomowa paczka: błędny punkt odrzuca całość; utrata odpowiedzi
po zapisie nie wycofuje danych.\\
FR-03 & Widoczność head poprzedza odpowiedź. None jest domyślne jak
w branchu; Written jawnie wymaga pełnego zapisu do OS. Brak gwarancji
per żądanie dotyczącej fsync; kryteria trwałości mierzy się w Written.\\

FR-04 & Zakres to obliczenia na siatce czasu. Selektor wybiera ostatnią
próbkę; agregacja wymaga funkcji.\\
FR-08 & Czas buforowania biegnie od przyjęcia najstarszej próbki
generacji. Flush całej generacji może wcześniej usunąć młodsze dane
z RAM; limit pamięci może go przyspieszyć.\\
FR-09 & Out-of-order wyłącznie w otwartym buforze: obowiązują
jednocześnie okno opóźnienia i granica zamknięcia.\\
NFR-03 & RSS $\leq1$ GiB (nie Gb); 288 MB surowych próbek nie obejmuje
zapasowej pojemności, indeksów ani zapytań.\\
NFR-04 & Brak utraty przy awarii procesu dotyczy Written, przy sprawnym
OS i nośniku. None dopuszcza utratę i jest mierzone osobno.
Błąd I/O WAL kończy proces; awaria hosta/nośnika jest poza gwarancją API.\\

\end{longtable}

\subsection{Rejestr głównych decyzji}
\begin{longtable}{P{13mm}P{143mm}}
\tabhead{ID & Decyzja, uzasadnienie i koszt}
DA-01 & Modularny monolit Rust: wywołania wewnątrz procesu upraszczają
transakcje i diagnostykę; skalowanie do jednego węzła.\\
DA-02 & Jeden Index, osobna blokada każdego Series: równoległe zapisy
rozłącznych szeregów, krótkie blokady katalogu, bez shardów.\\
DA-03 & Posortowane, niekompresowane bloki aktywne i zamrożone:
out-of-order i dostępność odczytu podczas flush.\\
DA-04 & Generacje head i jeden flush w tle: bariera zapisu obejmuje
drain i sync przez cut(), potem freeze. Zapis segmentu/manifestu
odbywa się po jej zwolnieniu.\\

DA-05 & Grafana jako jedyne GUI: standardowe źródło Prometheus,
zgodność sprawdzana z przypiętą wersją.\\
DA-06 & Zamknięte zakresy są niezmienne: proste indeksy i rozłączne
bloki szeregu; wcześniejszy flush może spowodować odrzucenie punktu.\\
DA-07 & WAL z 23-wal-improvements: bufor RAM, flusher/syncer, watch
Written, pliki bez nagłówków i checkpoint po cut()/remove\_through().\\
DA-08 & Atomowa paczka i spójny widok wybranych Series: blokady rosnąco po ID;
duża operacja może ograniczyć równoległość.\\
DA-09 & None domyślnie, Written per żądanie. Syncer działa okresowo,
bez publicznego Synced; LSN jest lokalny dla otwarcia WAL-a.\\

\end{longtable}

\subsection{Podstawa implementacyjna i zakres dostosowania}
\label{sec:wal-baseline}
Punktem odniesienia jest lokalny branch \code{23-wal-improvements},
commit \code{e1354b9} (,,docs'') \cite{wal23}.
Zachowujemy jego Wal/WalWriter, format kodeka, None/Written, okresowy syncer,
cut i ReplaySink. Pozostała część dokumentu opisuje docelowe spięcie
tego WAL-a z równoległymi Series i blokami aktywnymi/zamrożonymi;
nie przenosi globalnej blokady Db ani obecnego HTTP z brancha.

Wymagane korekty integracyjne zestawia sekcja~\ref{sec:wal-corrections}.
Są projektem do wykonania, a nie informacją, że już istnieją w branchu.
W szczególności nie dodaje się trybu Synced, globalnego trwałego LSN
ani innego formatu rekordów WAL.

\section{Architektura systemu}
\label{sec:architecture}

\subsection{Przypadki użycia i aktorzy}
Producent wysyła pojedyncze pomiary lub atomowe paczki. Analityk
wybiera czas i serwis w Grafanie albo odpytuje API. Operator uruchamia,
konfiguruje i restartuje bazę oraz kontroluje jej stan.

\begin{figure}[htbp]
\centering
\begin{tikzpicture}
\draw[draw=accent,rounded corners] (0.6,0.6) rectangle (11,-6.3);
\node[font=\small\bfseries,anchor=north] at (5.8,0.5)
  {TSDB z interfejsem Grafany};
\foreach \x/\y/\id/\name in {-1/-.9/prod/Producent,-1/-4.2/anal/Analityk,
  12.4/-2.6/ops/Operator} {
  \coordinate (\id) at (\x,\y);
  \draw[thick] (\x,\y+.6) circle (.15);
  \draw[thick] (\x,\y+.45) -- (\x,\y-.12)
    (\x-.33,\y+.25) -- (\x+.33,\y+.25)
    (\x,\y-.12) -- (\x-.28,\y-.5)
    (\x,\y-.12) -- (\x+.28,\y-.5);
  \node[font=\small,below] at (\x,\y-.55) {\name};
}
\node[usecase] (ingest) at (3.6,-1.2) {Zapis pomiarów\\FR-01, 02, 03, 09};
\node[usecase] (query) at (3.6,-3.4) {Zapytanie chwilowe\\lub zakresowe\\FR-04, 05, 06};
\node[usecase] (dash) at (3.6,-5.4) {Analiza dashboardu\\FR-07};
\node[usecase,text width=28mm] (run) at (8.6,-1.5)
  {Konfiguracja, start\\i restart bazy};
\node[usecase,text width=28mm] (status) at (8.6,-4)
  {Kontrola zdrowia\\i gotowości};
\draw[assoc] (prod) -- (ingest);
\draw[assoc] (anal) -- (query);
\draw[assoc] (anal) -- (dash);
\draw[assoc] (ops) -- (run);
\draw[assoc] (ops) -- (status);
\end{tikzpicture}
\caption{Diagram przypadków użycia UML. Flush (FR-08) jest automatycznym
mechanizmem silnika, a nie osobnym działaniem użytkownika.}
\label{fig:usecases}
\end{figure}

\subsection{Wdrożenie i granice procesów}
System działa na jednym komputerze. Przeglądarka łączy się z Grafaną;
serwer Grafany odpytuje TSDB, a generator wysyła HTTP JSON.
W środowisku nie ma serwera Prometheus.

\begin{figure}[htbp]
\centering
\begin{tikzpicture}[font=\small]
\draw[draw=accent,thick,fill=black!1] (0,0) rectangle (14,-9.1);
\node[anchor=north west,align=left] at (.2,-.15)
  {\textit{<<device>>} Host Linux x86-64};
\node[external,text width=40mm] (browser) at (3,-1.45)
  {Przeglądarka\\Grafana UI};
\node[external,text width=40mm] (curl) at (10,-1.45)
  {Klient testowy\\HTTP / curl};
\draw[draw=black!50,dashed] (.5,-2.8) rectangle (13.5,-6.9);
\node[anchor=north,font=\scriptsize,align=center] at (7,-2.95)
  {\textit{<<executionEnvironment>>}\\Docker Compose, sieć tsdb-net};
\node[box,text width=38mm] (grafana) at (3,-4.4)
  {\textit{<<container>>}\\grafana :3000};
\node[box,text width=42mm] (db) at (10,-4.4)
  {\textit{<<container>>}\\tsdb / Rust :8080};
\node[box,text width=38mm] (gen) at (3,-6.1)
  {\textit{<<container>>}\\generator};
\node[external,text width=42mm] (config) at (3,-8.2)
  {\textit{<<artifact>>}\\YAML źródła + JSON dashboardu};
\node[external,text width=42mm] (data) at (10,-8.2)
  {Wolumen tsdb-data\\\code{/var/lib/tsdb}};
\draw[assoc] (browser) -- node[note,right]{HTTP\\127.0.0.1:3000} (grafana);
\draw[assoc] (curl) -- node[note,right]{HTTP\\127.0.0.1:8080} (db);
\draw[assoc] (grafana) -- node[note,above]{HTTP / JSON\\http://tsdb:8080} (db);
\draw[assoc] (gen.east) -- node[note,above]{HTTP JSON: POST /api/v1/write\_json}
  (10,-6.1) -- (db.south);
\draw[assoc] (db.south east) -- (11.2,-5.4) --
  node[note,right]{lokalne I/O} (11.2,-7.5) -- (data.north east);
\draw[assoc] (config.west) -- (.7,-8.2) |-
  node[note,pos=.6,left]{mount\\read-only} (grafana.west);
\end{tikzpicture}
\caption{Diagram wdrożenia UML z zasobami lokalnymi.
Porty hosta są publikowane wyłącznie na loopback; kontenery używają DNS tsdb.
Szczegóły dostępu określa sekcja~\ref{sec:network}.}
\label{fig:deployment}
\end{figure}

Jeden proces ma wyłączny dostęp do katalogu: systemowa blokada pliku
\code{LOCK} trwa tak długo jak jego otwarty deskryptor; pozostały po
awarii plik nie oznacza aktywnej blokady. Sieciowe systemy plików
nie są platformą docelową.

\subsection{Warstwy, moduły i zależności}
Warstwa HTTP odpowiada za protokół, warstwa aplikacyjna za operacje
zapisu i zapytań, a silnik za dane i pliki. Parser zależy wyłącznie od
typów języka. Silnik nie zna kodów HTTP, Grafany ani składni PromQL.

\begin{figure}[htbp]
\centering
\begin{tikzpicture}
\node[box,text width=95mm] (http) at (0,0)
  {Adapter HTTP: zapis, API Prometheus, stan procesu};
\node[box] (write) at (-3,-2) {Koordynator zapisu\\walidacja paczek};
\node[box] (query) at (3,-2) {Parser i wykonawca\\zapytań};
\node[box,text width=108mm] (engine) at (0,-4.2)
  {Silnik: Index, indeks etykiet, Series z własnymi blokadami i blokami};
\node[box] (wal) at (-3,-6.2) {Bufor WAL\\flusher i syncer};
\node[box] (segment) at (3,-6.2) {Flush w tle\\segmenty i manifest};
\draw[arrow] (http.south) -- ++(0,-0.4) -| (write.north);
\draw[arrow] (http.south) -- ++(0,-0.4) -| (query.north);
\draw[arrow] (write) -- (engine.north -| write);
\draw[arrow] (query) -- (engine.north -| query);
\draw[arrow] (engine.south -| wal) -- (wal);
\draw[arrow] (engine.south -| segment) -- (segment);
\end{tikzpicture}
\caption{Zależności logiczne. Konfiguracja i diagnostyka są współdzielonymi usługami pomocniczymi.}
\end{figure}

\begin{longtable}{P{33mm}P{123mm}}
\tabhead{Moduł & Odpowiedzialność, interfejs i zależności}
HTTP & Dekodowanie, limity, błędy i odpowiedzi; Axum/Serde,
wywołania zapisu i zapytań.\\
Model i walidacja & Wspólne typy, kanonizacja klucza, limity nazw/etykiet
i wartości; bez HTTP i plików.\\
Koordynator zapisu & Pula, kolejność blokad, walidacja/rezerwacja;
Index, Series, WAL i bariera head.\\
Index i indeks & Klucz $\leftrightarrow$ Series, rejestracja i indeks etykiet;
metadane pod krótką blokadą katalogu.\\
Series i bloki & Posortowane próbki, duplikaty, granice czasu, aktywny/zamrożony
blok; blokada per Series, globalny budżet RAM.\\
WAL & Bufor, lokalny LSN, Written, okresowy sync i rotacja/cut;
mutex, flusher/syncer, watch, bez blokad Series w I/O.\\
Magazyn segmentów & Lokalne pliki, CRC, indeksy, manifest/checkpoint;
koordynator flush.\\
Parser & Tekst $\rightarrow$ AST/błąd z pozycją; typy języka,
bez dostępu do danych.\\
Wykonawca zapytań & Dobór Series, plan, agregacje i budżety;
AST i interfejs odczytu, bez dostępu do WAL.\\
Uruchomienie i diagnostyka & Konfiguracja, blokada katalogu, recovery,
gotowość, logi i zakończenie; zarządza cyklem życia modułów.\\
\end{longtable}

\subsection{Współbieżność i spójność}
\label{sec:concurrency}
Tokio obsługuje sieć; pula wykonuje pamięciową część zapisu
(domyślnie cztery zadania). Każdy Series ma własny RwLock.
Katalog zwalnia się przed pobraniem blokad Series.

\begin{longtable}{P{25mm}P{131mm}}
\tabhead{Operacja & Blokady i punkt publikacji}
Zapis & Dla nowych kluczy: mutex rejestracji; potem współdzielona
bariera generacji i wyłączne blokady Series rosnąco po ID.
Walidacja/rezerwacja $\rightarrow$ WalWriter.append $\rightarrow$
publikacja całej paczki przed zwolnieniem którejkolwiek blokady Series.\\
Append WAL & Thread-local scratch koduje poza mutexem.
Mutex bufora chroni kopię ramki i przyrost LSN; brak I/O, head append
i wait pod nim. Możliwe oczekiwanie na mutex/budżet, bez lock-free.\\
Written & Po publikacji zwolnić blokady, barierę i worker;
HTTP asynchronicznie wywołuje WalWriter.wait. None odpowiada bez wait.
Wątki WAL nie pobierają blokad Series.\\
ReadView & Pobrać uchwyty i wersję Index, zwolnić katalog,
pobrać współdzielone blokady wybranych Series rosnąco po ID
(try\_read z timeoutem). Zmiana wersji katalogu: zwolnić i ponowić.
Skopiować active i uchwyty frozen/dysku, zwolnić Series.\\
Checkpoint & Wyłączna bariera czeka na enqueue--head append;
bez blokad Series wykonuje blokujące Wal.cut, następnie zamraża head
i zmienia generację. Sync w cut może wstrzymać zapisy.
Odczyty nie pobierają bariery; segment i manifest powstają po jej
zwolnieniu (sekcja~\ref{sec:checkpoint}).\\
\end{longtable}

Jeden ReadView obejmuje całe query\_range, odpowiada chwili stanu
wybranych Series i nie zawiera części paczki. Może zobaczyć dane klienta
czekającego na Written. Publikacje rozłącznych Series mogą kończyć się
poza kolejnością WAL; wspólne Series zachowują tę kolejność.
Odczyt po sukcesie widzi zapis w swoim zakresie.
Append sam nie wykonuje I/O, ale żądanie może czekać na barierę cut.

\section{Model danych i kontrakty modułów}
\label{sec:model}

Nazwy odpowiadają kodowi: \code{Index} (\code{src/storage/index.rs}),
\code{Series} i \code{SeriesState} (\code{src/storage/series.rs})
oraz \code{LabelSet}, \code{Sample}, \code{SeriesId} i \code{TimeRange}
(\code{src/model.rs}). Bazą zarządza \code{Db} z \code{DbState}.
Metody to \code{Db.write(WriteBatch)}, \code{Db.wait(Lsn)}
i \code{Db.query(&[Matcher], TimeRange)} zwracająca \code{Vec<SeriesResult>}.
Pliki danych obsługuje \code{SegmentStore} (\code{src/storage/segment.rs}).
Diagram pokazuje istniejące pola i metody oraz planowane bloki.
\code{MutableBlock}, \code{FrozenBlock}, \code{BlockDescriptor},
\code{ReadView}, granica \code{sealed_through}, \code{last_lsn}
i wersjonowanie katalogu są elementami projektu do implementacji.
Obecny \code{Series.state} używa \code{Mutex<SeriesState>}; opisane dalej
współdzielone blokady odczytu wymagają docelowo \code{RwLock}.


\subsection{Tożsamość szeregu i reprezentacja punktu}
Klucz szeregu składa się z nazwy metryki oraz uporządkowanej listy par
etykiet. Klucze etykiet są unikalne; sortowanie jest leksykograficzne
według bajtów UTF-8. Nie normalizuje się wielkości liter ani Unicode.
Porównanie pełnego klucza rozstrzyga tożsamość również przy kolizji
funkcji skrótu.

\begin{longtable}{P{32mm}P{46mm}P{78mm}}
\tabhead{Typ / pole & Reprezentacja & Reguła}
\code{LabelSet} & Metryka i posortowane pary napisów. &
Nazwa metryki: \code{[a-zA-Z_:][a-zA-Z0-9_:]*}, do 128 bajtów.\\
\code{Label} & Nazwa i wartość UTF-8. &
Nazwa: \code{[a-zA-Z_][a-zA-Z0-9_]*}, do 64 bajtów;
wartość do 256 bajtów, maksymalnie 16 etykiet.\\
Nazwy zastrzeżone & Prefiks \code{__}. &
Nie są dopuszczane w etykietach zapisu. \code{__name__} jest wirtualną
etykietą nazwy metryki w odpowiedziach.\\
\code{SeriesId} & \code{u64}, od 0. &
Rosnący identyfikator bieżącego head; po restarcie może być inny.
Występuje tylko w head i WAL, nie w API, manifeście ani segmentach.\\
\code{Sample.timestamp} & \code{u64}, milisekundy od epoki Unix. &
Zakres przyjmowania: $0$ do $2^{53}-1$ ms, aby zachować dokładność
całkowitą w klientach JSON.\\
\code{Sample} & \code{timestamp: u64}, \code{value: f64}. &
16 bajtów w bloku; wartość wejściowa musi być skończona.
Normalizacja $-0$ do $+0$.\\
\code{Lsn} & \code{u64}, lokalny dla Wal::open. &
Liczba bajtów ramek przyjętych w tym otwarciu WAL-a; 0 przed pierwszą.
Nie przechodzi przez restart i nie jest trwałym kursorem recovery.\\

\end{longtable}

Wewnętrzny \code{LabelSet} zawiera także \code{__name__} z nazwą metryki.
Adapter tworzy go z metryki i etykiet; nazwa zastrzeżona nadal nie jest dozwolona
w wejściowym obiekcie labels. WalWriter używa LabelSetEncoding z brancha.

Pusta wartość etykiety jest kanonizowana przez usunięcie tej etykiety.
Dzięki temu porównanie z pustym napisem może objąć również szeregi
bez danej etykiety. Powtórzona nazwa klucza w obiekcie JSON jest błędem
przed kanonizacją, także gdy parser JSON domyślnie potrafiłby ją nadpisać.
Nieznane pola wejściowe są odrzucane. Wartości metryk i etykiet
w przykładach stosują konwencje Prometheusa; projekt celowo ogranicza
dozwolone nazwy do wskazanej składni ASCII.

\subsection{Struktury Index, Series i bloki}
\label{sec:tsmap}
Inspiracją jest oddzielenie katalogu, synchronizacji szeregu i bloków
otwartych/zamkniętych w Gorilla, sekcje 4.2--4.3 \cite{gorilla}.
Tutaj jeden Index obsługuje całą bazę, a bloki bez kompresji zapewniają
dostępność podczas flush. Buforowany WAL ma jawne poziomy trwałości.

\begin{figure}[htbp]
\centering
\begin{tikzpicture}
\node[umlclass,text width=56mm] (map) at (0,0) {
  \textbf{Index}\nodepart{second}
  inverted: HashMap<LabelSet, Arc<Series>{}>\\
  forward: HashMap<SeriesId, Arc<Series>{}>\\
  posting\_index: HashMap<String,\\
  HashMap<String, Vec<SeriesId>{}>{}>\\
  all\_ids: Vec<SeriesId>, next\_id: SeriesId
  \nodepart{third}encode(labels), get(id)\\resolve(matchers), id\_for(labels)};
\node[umlclass,text width=45mm] (ts) at (8,0) {
  \textbf{Series}\nodepart{second}
  id: SeriesId\\labels: Arc<LabelSet>\\
  state: Mutex<SeriesState>
  \nodepart{third}new(id, labels), lock()};
\node[umlclass,text width=50mm] (st) at (8,-3.4) {
  \textbf{SeriesState}\nodepart{second}
  samples: Vec<Sample>\\
  newest: Option<u64>\\logged\_gen: u64
  \nodepart{third}accepts(timestamp), append(sample)\\
  range(TimeRange), take\_samples()\\
  is\_logged(g), mark\_logged(g)};
\node[umlclass,text width=43mm] (mut) at (0,-3.4) {
  \textbf{MutableBlock} (plan)\nodepart{second}
  samples: Vec<Sample> \{sorted\}
  \nodepart{third}insert(sample), slice(bounds)};
\node[umlclass,text width=43mm] (froz) at (0,-6.3) {
  \textbf{FrozenBlock} (plan)\nodepart{second}
  samples: Vec<Sample>\\\{immutable, shared by Arc\}
  \nodepart{third}slice(bounds)};
\node[umlclass,text width=50mm] (disk) at (8,-6.3) {
  \textbf{BlockDescriptor} (plan)\nodepart{second}
  segment, min/max, offset, count
  \nodepart{third}read(bounds)};
\draw[assoc] (map) -- node[note,pos=.12,above]{1}
  node[note,pos=.83,above]{0..*} (ts);
\draw[compose] (ts) -- node[note,right]{1 : 1} (st);
\draw[compose] (st) -- node[note,above]{active: 1} (mut);
\draw[assoc] (st.south west) -- ++(-0.8,0) |- node[note,pos=.7,above]
  {frozen: 0..1} (froz.east);
\draw[compose] (st) -- node[note,right]{persisted: 0..*} (disk);
\end{tikzpicture}
\caption{Diagram klas modelu head (typy Rust, bez dziedziczenia).
Romb oznacza własność; zwykłe linie także współdzielone uchwyty.
Klucz: LabelSet, czas: u64; Sample = timestamp: u64 + value: f64 (16 B).}
\label{fig:classes}
\end{figure}

W kodzie \code{Db.state: RwLock<DbState>} chroni \code{Index}.
\code{inverted} mapuje LabelSet na szereg, \code{forward} mapuje SeriesId
na szereg, \code{posting_index} zawiera listy ID według etykiet,
a \code{all_ids} służy do skanowania katalogu. ID nie jest indeksem wektora. Szeregów nie usuwa się, a ID nie wykorzystuje ponownie.

Rejestracja nowych Series wymaga mutexu i ponownego lookup. Przygotowuje je
poza publiczną mapą; po udanym append i inicjalizacji head krótka
blokada katalogu publikuje Series/indeksy i zwiększa wersję.
Katalog nie jest blokowany przez kodowanie WAL ani I/O;
nieudana paczka może pozostawić lukę ID, bez pustego publicznego Series.

Każdy Series przechowuje \code{logged_gen}. Pierwszy zapis w generacji $g$
dołącza do ramki jego definicję, nawet jeśli Series istniał przed cut.
Znacznik ustawia się dopiero po append, pod tą samą blokadą Series.
Nie tylko create callback decyduje o definicji.
Dzięki temu pliki po checkpointcie nie zależą od definicji usuniętych
przez remove\_through. Zwykła rotacja pliku nie zmienia generacji.
Po restarcie Series dostają nowe ID i nieustawiony logged\_gen;
GC head pozostaje poza MVP.

\subsection{Duplikaty i atomowość paczek}
W szeregu jest najwyżej jedna wartość dla czasu. Identyczny punkt
w active to duplikat bez zmian; inna wartość odrzuca całą paczkę,
również dla powtórzeń w niej. Zamknięty zakres odrzuca także identyczny
duplikat bez odczytu dysku. Brak bezterminowej idempotencji/rejestru kluczy;
po utracie odpowiedzi wynik można sprawdzić zapytaniem.

Walidacja tworzy plan całej paczki; jej definicje i próbki trafiają do
jednej ramki. Samych duplikatów nie loguje się. Ich Written czeka na
maksimum last\_lsn dotkniętych Series, zamiast na globalny koniec bufora
zwracany przez pusty append biblioteki. Historyczne dane po otwarciu
zaczynają z last\_lsn=0, nowe append ustawiają lokalny LSN.
Nie jest to obietnica widoczności mniejszych LSN w rozłącznych Series.

\subsection{Kontrakty wewnętrzne}
Notacja opisuje granice modułów bez narzucania crate'ów.
Poza istniejącym API Db i WAL poniższe kontrakty (m.in. submit,
snapshot i evaluate) są propozycjami do implementacji.
Operacje mogą zawieść; wyjątkiem w branchu jest wait, którego
domknięcie błędem wymaga korekty W-04 z sekcji~\ref{sec:wal-corrections}.
Selektory/AST/widoki są pożyczane, QueryBudget mutowalny.

\begin{longtable}{P{72mm}P{84mm}}
\tabhead{Operacja $\rightarrow$ wynik & Kontrakt}
\code{canonicalize(Measurement)} $\rightarrow$ \code{CanonicalPoint};
\code{parse(&str)} $\rightarrow$ \code{Expr} &
LabelSet i walidacja; AST selektora, funkcji albo skalara.\\
\code{submit(Vec<CanonicalPoint>, Durability)} $\rightarrow$
\code{WriteReceipt} & Widoczna cała paczka i osiągnięty tryb.
Receipt: inserted, duplicates, lsn, durability.\\
\code{snapshot(Selector, TimeRange, QueryBudget)} $\rightarrow$
\code{ReadView}; \code{read_series(ReadView, SeriesId, TimeRange)}
$\rightarrow$ \code{SampleCursor} &
Wersja katalogu, kopie active, uchwyty frozen/dysku;
próbki chronologiczne, bez ujawniania źródła.\\
\code{evaluate(Expr, EvalGrid, ReadView, QueryBudget)} $\rightarrow$
\code{QueryResult} & Budżet czasu, RAM, próbek i wyniku.\\
\code{WalWriter.append(series, samples)} $\rightarrow$ \code{Lsn};
\code{append_samples(samples)} $\rightarrow$ \code{Lsn} &
Iteratory definicji (ID, LabelSet) i grup (ID, \code{&[Sample]}).
Kodowanie i kopia do RAM bez I/O; błąd WalError.\\
\code{WalWriter.wait(Lsn, Durability)}; \code{written_lsn()};
\code{backlog_bytes()} & None/Written, watch i pobudzenie flushera;
czekanie dopiero po publikacji head.\\
\code{Wal.cut()} $\rightarrow$ \code{u64};
\code{remove_through(u64)} $\rightarrow$ \code{usize} &
Cut drenuje, synchronizuje i rotuje; zwraca ID zamkniętego pliku.
Usuwanie tylko po trwałym checkpointcie, nigdy active.\\
\code{freeze_head()} $\rightarrow$ \code{FrozenGeneration};
\code{persist(FrozenGeneration)} $\rightarrow$ \code{Checkpoint} &
Frozen zawiera $g$, wal\_checkpoint=$N$, klucze/granice i bloki.
Persist poza barierą.\\
\code{Wal.open(WalConfig, ReplaySink)}; \code{Wal.close()} &
Odtworzenie przez sink i nowy plik; close drenuje/synchronizuje
i dołącza wątki. Filtrowanie po $N$: korekta W-01.\\
\end{longtable}

Błędy modelu/parsera/zapisu/odczytu/zapytań są typowane.
WalError zachowuje Io, RecordTooLarge, Corrupt i Closed;
przeciążenie jest błędem przyjęcia przed append.
Po append publikacja korzysta z rezerwacji, bez zwykłego rollbacku;
timeout/anulowanie nie wycofują danych. I/O WAL kończy proces,
pozostałe błędy mapuje sekcja~\ref{sec:http-errors}.

\subsection{Podział odpowiedzialności zespołu}
\begin{tabularx}{\linewidth}{P{43mm}Y}
\toprule
Osoba & Obszar odpowiedzialności\\ \midrule
Jan Zaborski & WAL, segmenty, manifest, indeksy dyskowe, checkpoint i recovery.\\
Hubert Błonowski & Parser, wykonawca zapytań, adapter HTTP, integracja i dashboard Grafany.\\
Obaj autorzy & Model danych, bufor i granice jego interfejsu, przegląd kontraktów,
testy integracyjne, pomiary i dokumentacja swoich podsystemów.\\
\bottomrule
\end{tabularx}

Wykonawcę zapytań można rozwijać z pamięciową implementacją interfejsu
odczytu. Magazyn można testować bez HTTP i parsera. Adapter HTTP korzysta
z ustalonych typów wyników i błędów. Wspólna akceptacja zmian modelu
i kontraktów poprzedza ich wdrażanie po obu stronach.

\section{Zapis, buforowanie i obsługa opóźnionych danych}
\label{sec:write}


\subsection{Algorytm przyjęcia paczki}
\begin{enumerate}
\item Przed wczytaniem ciała zarezerwować przyjęcie;
sprawdzić rozmiar/tryb, zdekodować, kanonizować i pogrupować po Series.
\item Pobrać blokady z sekcji~\ref{sec:concurrency}. Zweryfikować
całą paczkę oraz zarezerwować head, indeksy, scratch i miejsce WAL.
Brak miejsca: zwolnić blokady/barierę, czekać i walidować ponownie.
Paczka niemieszcząca się w pustym budżecie daje 422; przejściowy brak
po maks. 1 s daje 503, zawsze przed append.
\item Zebrać definicje Series z logged\_gen innym niż $g$ i nowe próbki.
WalWriter.append koduje jedną ramkę w scratch i kopiuje ją do WalBuffer.
\item Pod blokadami Series opublikować całość, last\_lsn i logged\_gen=$g$.
Opublikować nowe Series/indeksy; zwolnić blokady, barierę i worker.
\item None odpowiada; Written czeka asynchronicznie na written\_lsn.
\end{enumerate}
Recovery może odtworzyć pełną ramkę bez odpowiedzi lub zakończonego
head append. Publikacja i synchronizacja nie tworzą jednej transakcji.

\begin{figure}[htbp]
\centering
\begin{tikzpicture}
\foreach \x/\name in {0/HTTP,3/Worker,6/WalBuffer,9/Flusher,12/Syncer} {
 \node[font=\small\bfseries] at (\x,.3) {\name};
 \draw[dashed,black!35] (\x,0) -- (\x,-5.4);
}
\draw[arrow] (0,-.5) -- node[note,above]{paczka} (3,-.5);
\draw[arrow] (3,-1.2) -- node[note,above]{append} (6,-1.2);
\draw[arrow,dashed] (6,-1.8) -- node[note,above]{lokalny LSN} (3,-1.8);
\node[note] at (3,-2.45) {publikacja head\\zwolnienie blokad};
\draw[arrow,dashed] (3,-3.1) -- node[note,above]{None gotowe} (0,-3.1);
\draw[arrow] (6,-3.6) -- node[note,above]{swap / write\_all} (9,-3.6);
\draw[arrow,dashed] (9,-4.35) -- node[note,above]{watch Written} (0,-4.35);
\node[note] at (12,-4.9) {okresowe fdatasync\\bez waitera klienta};
\end{tikzpicture}
\caption{Ścieżka zapisu zgodna z WalWriter. None kończy po head append,
Written po write\_all; okresowy syncer nie potwierdza trybu per żądanie.}
\end{figure}

\subsection{Warunki przyjmowania out-of-order}
$H_s$ to największy przyjęty czas szeregu (także na dysku), $S_s$ --
największy czas danych zamrożonych/opublikowanych, a $W$ -- okno
opóźnienia. Dla każdego nowego punktu:
\[
 t>S_s \quad\land\quad t\geq H_s-W.
\]
$S_s$ pomija się przed pierwszym zamrożeniem, $H_s$ dla nowego Series.
Całą paczkę waliduje się względem $H_s$ sprzed paczki, niezależnie od
kolejności JSON; po publikacji $H_s$ jest maksimum starego i nowych czasów.
Arytmetyka jest sprawdzana, ujemna dolna granica oznacza zero.

Kolejność decyzji: zamknięty zakres $\rightarrow$ duplikat/konflikt
w otwartym buforze $\rightarrow$ test $W$ dla nowego punktu.
Identyczny otwarty duplikat nie staje się nowym zapisem.
Próbka może wyprzedzać zegar serwera najwyżej o 60 s, co chroni $H_s$
przed błędnym przesunięciem. Zegar ścienny służy temu testowi i domyślnemu
czasowi zapytania; pobyt w RAM i timeouty używają zegara monotonicznego.

\begin{tabularx}{\linewidth}{P{63mm}P{28mm}Y}
\toprule
Stan & Czas punktu & Wynik\\ \midrule
$H=12{:}10$, $S=12{:}00$, $W=5$ min & 12:06 / 12:04 &
Przyjęcie / przekroczone $W$.\\
Po zamrożeniu $S=12{:}10$ & 12:09 / 12:11 &
Zakres zamknięty / nowy aktywny blok (także podczas flush).\\
\bottomrule
\end{tabularx}

Przy $W=0$ paczka może zawierać kilka czasów, ale żaden nowy nie może
poprzedzać starego $H_s$; punkty paczki nie ustalają kolejności przyjęcia.
Profil standardowy: $W=5$ min.

\subsection{Organizacja i limit pamięci}
Posortowany wektor dopisuje rosnące czasy na końcu; out-of-order wymaga
wyszukiwania binarnego i liniowego przesunięcia. Dwa wyszukiwania
wyznaczają wycinek odczytu. Pojemność rośnie przez
\code{try_reserve_exact} w wielokrotnościach 4096 próbek;
rozlicza się faktyczną pojemność, także zapas i przejściowe
współistnienie starej/nowej alokacji.

Budżet head obejmuje aktywne, zamrożone i podtrzymywane przez ReadView
bloki, każdy współdzielony blok raz do ostatniego uchwytu.
Katalog, przyjęcia, WAL i zapytania mają osobne budżety.
Rezerwacja poprzedza enqueue; błąd nie publikuje próbek.
Za mały limit minimalnej alokacji oznacza błąd konfiguracji/paczki,
bez nieskończonego flush.

Zamrożenie przenosi wektor do \code{Arc<FrozenBlock>} bez pełnej kopii.
Nowy aktywny blok alokuje w miarę napływu danych. Zapytania kopiują
wyłącznie potrzebne aktywne próbki; niezmienne dane współdzielą.
Po publikacji segmentu starszy widok może nadal podtrzymywać RAM,
co uwzględnia się w budżecie i pomiarach RSS.

\subsection{Wyzwalanie i przebieg zrzutu}
Timer co 1 s zleca flush po upływie czasu pobytu najstarszej próbki
generacji w RAM; termin dotyczy rozpoczęcia, nie zwolnienia pamięci.
Inne wyzwalacze: przekroczenie 80\% budżetu lub brak miejsca na paczkę.
Pustej generacji nie zrzuca się.

Koordynator wykonuje cut WAL i zamraża razem niepuste bloki.
Po zwolnieniu bariery nowe zapisy działają równolegle z jednym flush
w tle (sekcja~\ref{sec:checkpoint}). Do jego zakończenia następny flush
pozostaje oczekujący, bez kolejki zamrożonych generacji.
Przy twardym limicie paczki krótko czekają bez blokad lub dostają 503;
dostępność podczas flush podlega ograniczonej pamięci i szybkości dysku.

\begin{figure}[htbp]
\centering
\begin{tikzpicture}
\node[state] (active) at (0,0) {Aktywny\\generacja g};
\node[state] (frozen) at (4.8,0) {Zamrożony\\generacja g};
\node[state] (disk) at (9.6,0) {Na dysku\\generacja g};
\node[state] (next) at (0,-2.4) {Nowy aktywny\\generacja g+1};
\node[state] (io) at (6.5,-2.4) {Flush w tle\\odczyt g dostępny};
\draw[arrow] (active) -- node[note,above=2mm]{cut + freeze} (frozen);
\draw[arrow] (frozen) -- node[note,above=2mm]{trwały checkpoint} (disk);
\draw[arrow] (active) -- node[note,left]{nowe zapisy} (next);
\draw[arrow] (frozen) -- (io);
\draw[draw=black!50,dashed] (next.east) --
  node[note,above]{działają równolegle} (io.west);
\end{tikzpicture}
\caption{Cut i freeze tworzą nowy active przed zapisem frozen do segmentu.
Kolejna rotacja czeka na zakończenie poprzedniego flush.}
\end{figure}

\section{Organizacja danych na dysku}
\label{sec:storage}

\subsection{Katalog i wspólne reguły formatu}
\begin{lstlisting}
data/LOCK
data/MANIFEST
data/wal/wal-00000001.log
data/wal/wal-00000002.log
data/segments/segment-00000000000000000001.bin
\end{lstlisting}
Nazwy WAL mają minimum 8 cyfr jak file\_name() w branchu, segmentów
danych 20. Ich numeracje są niezależne. Jeden WAL jest aktywny.
Pliki WAL nie mają nagłówków ani .tmp; create\_new i sync katalogu
publikują nowy, początkowo pusty plik. Segmenty/manifest powstają przez
.tmp, sync, rename i sync katalogu. Manifest wskazuje segmenty oraz
największy zamknięty WAL objęty checkpointem.

Liczby: little-endian, timestamp u64, f64 jako bity IEEE 754,
bez paddingu/układu struktur Rust. Napis: u32 długości i UTF-8.
Arytmetyka/limity są sprawdzane przed alokacją i odczytem.
CRC-32/IEEE \cite{crc}, zapis u32 LE, wektor
\code{123456789 -> 0xCBF43926}; wykrywanie uszkodzeń, bez uwierzytelnienia.

WAL zachowuje format brancha bez pola wersji. Docelowy indeksowany
segment ma wersję 3, manifest 4; ich formaty zastępują wersje z dokumentu
1.3, aby nie utrwalać head SeriesId ani lokalnych LSN.
Nie są formatem obecnego SegmentStore JSON z brancha.
Nieznana wersja, flagi, UTF-8 i długości zatrzymują start;
migracja plików prototypu wymaga osobnego narzędzia.

\subsection{Bufor WAL, flusher i syncer}
\label{sec:wal-runtime}
Usługa zachowuje \code{Wal}, klonowalny \code{WalWriter} i
\code{Shared} z brancha. \code{WalBuffer} chroni Vec, end i closed
jednym mutexem; flusher podmienia Vec z własnym spare.
LSN rośnie o pełną długość ramki w jednym otwarciu:
\[
 0\leq written\_lsn\leq buffer.end,\qquad
 backlog\_bytes=buffer.end-written\_lsn.
\]
Po ponownym Wal::open oba liczniki zaczynają od 0. Nie ma
synced\_lsn ani publicznego oczekiwania na sync.

\begin{tabularx}{\linewidth}{P{23mm}Y}
\toprule
Tryb & Warunek sukcesu po publikacji head\\ \midrule
None (domyślny) & Bez oczekiwania na I/O; awaria procesu może utracić
potwierdzoną paczkę. Zachowane pełne ramki mogą wrócić.\\
Written & \code{written_lsn() >= lsn}; pełna ramka w OS/page cache.
Potwierdzenie przetrwa awarię samego procesu przy sprawnym hoście/nośniku,
bez gwarancji dla awarii hosta.\\
\bottomrule
\end{tabularx}
\par


\begin{figure}[htbp]
\centering
\begin{tikzpicture}
\node[umlclass,text width=48mm] (writer) at (0,0) {
 \textbf{WalWriter}\nodepart{second}shared: Arc<Shared>
 \nodepart{third}append(series, samples)\\wait(lsn, None | Written)};
\node[umlclass,text width=48mm] (shared) at (8,0) {
 \textbf{Shared}\nodepart{second}
 buffer: WalBuffer\\progress: Progress (watch Written)\\
 sync\_file: Mutex<Arc<File>{}>\\active\_segment, cut
 \nodepart{third}progress.publish\_written()\\progress.unpark\_flusher()};
\node[umlclass,text width=48mm] (flush) at (0,-3.5) {
 \textbf{Threads::flush\_loop}\nodepart{second}spare: Vec<u8>\\seg: ActiveSegment
 \nodepart{third}buffer.take(), write\_batch()\\rotate(), cut.finish()};
\node[umlclass,text width=48mm] (sync) at (8,-3.5) {
 \textbf{Threads::sync\_loop}\nodepart{second}last: Lsn
 \nodepart{third}co sync\_interval:\\clone(sync\_file), sync\_data()};
\draw[assoc] (writer) -- (shared);
\draw[arrow,dashed] (flush.north east) -- node[note,sloped,above]
 {take / Written} (shared.south west);
\draw[arrow,dashed] (shared) -- node[note,right]{Arc<File>} (sync);
\draw[arrow,dashed] (writer) -- node[note,left]{pobudzenie} (flush);
\end{tikzpicture}
\caption{Typy i funkcje WAL-a. Flusher i syncer to role wątków Threads.
Syncer nie publikuje postępu do klienta; cut obsługuje flusher.}
\label{fig:wal-classes}
\end{figure}


\textbf{Append.}  Kodowanie w thread-local scratch, kopia pełnej ramki
pod mutexem, zwiększenie end; bez plikowego I/O.
Przekroczenie flush\_threshold budzi flusher. Pusty append biblioteki
zwraca bieżący end; aplikacja obsługuje no-op według
sekcji~\ref{sec:model}. Scratch większy od 1 MiB jest zmniejszany po użyciu.


\textbf{Flusher.}  Domyślnie budzony co 100 ms lub po 4 MiB,
także przez Written i cut/close. Poza cut/close przy waiterach i buforze mniejszym
od commit\_bytes=256 KiB zbiera grupę przez 200~$\mu$s.
Swap pod mutexem, write\_all poza nim; Written dopiero po zapisaniu
całej przejętej porcji. Krótki write obsługuje write\_all
\cite{linux-write}. Nowe append używają drugiego Vec.


\textbf{Syncer.}  Co 1 s odczytuje Written; przy postępie klonuje
Arc aktywnego pliku i wykonuje sync\_data. Nie jest pobudzany żądaniem
trwałości i nie wysyła potwierdzenia. Rotacja sama synchronizuje stary
plik; jego Arc pozostaje ważny podczas równoległego sync.


\textbf{Rotacja rozmiarowa.}  Przed zapisem przejętego bufora sprawdza,
czy plik już ma co najmniej 128 MiB; jeśli tak, sync starego,
create\_new następnego i sync wal/. Próg jest miękki: zapisuje cały
przejęty bufor, bez dzielenia ramek, i może przekroczyć 128 MiB
o rozmiar porcji. Sama rotacja nie zmienia generacji head.


\textbf{Cut.}  Serializowany handshake przez mutex/Condvar budzi flusher,
który drenuje bufor, synchronizuje/zamyka active i tworzy następny.
Zwraca numer zamkniętego pliku. Aplikacja na ten czas blokuje nowe
append barierą generacji, aby cut zakończył drain i dokładnie pokrywał
zamrażany head. Wywołanie jest blokujące, poza wątkiem Tokio.


\textbf{Wait i close.}  Written rejestruje waiter, subskrybuje ostatni
stan watch i budzi flusher; nie wymaga każdego zdarzenia
\cite{tokio-watch}. Obsługa HTTP nakłada timeout 30 s, bez rollbacku.
Close odrzuca kolejne append, drenuje, wykonuje końcowy sync
i dołącza wątki. Domknięcie kanału nie może udawać sukcesu (W-04).


\code{writeback=true} uruchamia na Linux
\code{sync_file_range(..., SYNC_FILE_RANGE_WRITE)} po zapisie porcji;
na pozostałych platformach to no-op. To optymalizacja writeback,
nie fdatasync ani gwarancja trwałości \cite{linux-range}.
Błąd write/sync/writeback/tworzenia WAL kończy proces (fail-stop).
Sukcesu nie otrzymuje nieosiągnięty cel Written; None mogło odpowiedzieć.
Okres 1 s nie ogranicza straty None bez znajomości backlogu i opóźnień I/O.

\subsection{Wymagane korekty i kontrakty integracyjne}
\label{sec:wal-corrections}
Poniższe zmiany są uzupełnieniem brancha, nie przebudową jego WAL-a.
\begin{longtable}{P{13mm}P{143mm}}
\tabhead{ID & Różnica względem e1354b9 i wymagane zachowanie}
W-01 & Dodać WalConfig.replay\_after (u64, domyślnie 0),
ustawiane wyłącznie z manifestu na $N$. Reader/open mają
pomijać pliki $\leq N$, także gdy usuwanie przerwała awaria.
Nowy plik ma ID większe od $N$ i wszystkich znalezionych plików.
Nie dodaje się trwałego bajtowego kursora.\\
W-02 & Zastąpić gałąź readera odcinającą dowolny błąd ostatniego pliku
ograniczeniem do niepełnego końca (Truncated).
Pełne błędne CRC/kind/payload lub luka plików po $N$ zatrzymują start.
Kanoniczność LabelSet sprawdzać przed normalizacją dekodera.
Format ramek pozostaje bez zmian.\\
W-03 & Bufor brancha rośnie bez twardego limitu. Dodać rezerwację przed
append i limit pojemności active+spare i realokacji (32 MiB);
scratch rozliczać w 64 MiB przyjęcia. Początkowe 8 MiB każdego Vec
nie jest limitem. Backpressure bez blokad Series.\\
W-04 & Progress.wait nie może ignorować błędu zamknięcia watch.
Docelowo \code{Result<(), WalError>}, z Closed przy zamknięciu kanału;
timeout/cancel adaptera nie wycofują head.
Zachować jeden watch Written, bez nowej granicy Synced.\\
W-05 & Integracja head stosuje logged\_gen i barierę obejmującą cut.
Definicja wraca przy pierwszym zapisie Series po cut/restart.
Nie wystarcza logowanie tylko nowych szeregów.
ReplaySink scala po LabelSet, a nie przez utrwalony head ID.\\
W-06 & Resolver zachowuje buforowanie late series, mapowanie ID
i statystyki. Db nie ogłasza gotowości przy dropped\_samples $>0$,
braku budżetu pending lub błędzie semantycznym sinka.
Czytnik biblioteczny może raportować osierocone próbki; produkcyjny
start nie może ich cicho utracić.\\
\end{longtable}

\subsection{Format dziennika i lokalny LSN}
\label{sec:wal}
Plik jest ciągiem ramek od offsetu 0, bez magic, wersji i base\_lsn.
Ramka dokładnie odpowiada RecordHeader/RecordEncoder brancha:
\begin{longtable}{P{18mm}P{23mm}P{32mm}P{83mm}}
\tabhead{Offset & Rozmiar & Pole & Znaczenie}
0 & 4 B & len: u32 & Długość samego payloadu; maksimum biblioteki 64 MiB.\\
4 & 4 B & crc: u32 & CRC(kind || len w LE || payload), w tej kolejności.\\
8 & 1 B & kind: u8 & Write=0.\\
9 & len B & payload & Dwa liczniki, definicje, grupy próbek.\\
\end{longtable}
Cała ramka ma $9+len$ B. Profil HTTP ogranicza ją dodatkowo do 8 MiB
przed append; kodek zachowuje limit payloadu 64 MiB.
Pustego rekordu (oba liczniki 0) WalWriter nie dopisuje.

Payload: \code{series_count: u32}, \code{samples_count: u32}
(liczba grup, nie punktów), następnie:
\begin{enumerate}
\item Definicje: \code{id: u64}, \code{label_count: u32} i LabelSet:
pary napisów (u32 długości + UTF-8), w tym \code{__name__}.
\item Grupy: \code{id: u64}, \code{point_count: u32},
próbki \code{timestamp: u64, value: f64} po 16 B.
Pusta grupa jest pomijana przez encoder.
\end{enumerate}
Aplikacja porządkuje definicje/grupy po ID i punkty po czasie, choć
WAL nie wymaga sortowania grup. Dekoder wykrywa nadmierne liczniki,
UTF-8, brakujące i nadmiarowe bajty; adapter dodatkowo waliduje etykiety,
czasy, skończoność i konflikty. Jednostka replay to pełna ramka,
bez osobnego rekordu commit.

Próbka istniejącego Series: payload $8+12+16=36$ B, ramka 45 B.
Pierwsza paczka API dwóch punktów nowego Series ma LabelSet 74 B,
definicję $8+74=82$ B, payload $8+82+12+32=134$ B i ramkę 143 B.
Jeśli jest pierwszą po otwarciu, LSN=143 i EOF=143.
Po rotacji LSN nadal rośnie; po restarcie zaczyna od 0 niezależnie od EOF
starszych plików. Dlatego LSN z odpowiedzi nie porównuje się między
uruchomieniami ani nie zapisuje w manifeście.

\subsection{Segment danych}
\label{sec:segment}
Docelowy segment przechowuje klucze, nie head SeriesId, tak jak
DiskSeries w branchu. Własny indeksowany format ułatwia odczyt wycinków:
nagłówek 48 B, słownik LabelSet, bloki 16-bajtowych próbek,
indeks i CRC całego pliku (bez końcowego pola CRC).
Jeden flush obejmuje niepuste frozen generacji.

\begin{longtable}{P{18mm}P{23mm}P{34mm}P{81mm}}
\tabhead{Offset & Rozmiar & Pole & Znaczenie}
0 & 8 B & magic & ASCII \code{RSTSEG03}.\\
8 / 10 / 12 & 2 / 2 / 4 B & version / flags / header\_len &
u16=3 / u16=0 / u32=48.\\
16 & 8 B & generation & Numer generacji head, u64.\\
24 & 8 B & wal\_checkpoint & ID zamkniętego WAL $N$ z cut, nie LSN.\\
32 & 8 B & series\_count & Liczba kluczy/bloków/wpisów indeksu.\\
40 & 8 B & index\_offset & Początek indeksu od początku pliku.\\
\end{longtable}
Słownik od offsetu 48 zawiera LabelSetEncoding dla każdego bloku,
w kolejności kanonicznych kluczy; indeks wskazuje początek danego klucza.
Po słowniku następują w tej samej kolejności niepuste bloki próbek
\code{timestamp: u64, value_bits: u64}, czasy ściśle rosnące.
Indeks ma po 48 B na blok:
\begin{longtable}{P{18mm}P{23mm}P{34mm}P{81mm}}
\tabhead{Offset wpisu & Rozmiar & Pole & Znaczenie}
0 & 8 B & key\_offset & Offset LabelSet w słowniku; nie identyfikator Series.\\
8 / 16 & 8 / 8 B & min\_time / max\_time & Pierwszy/ostatni czas, u64.\\
24 / 32 & 8 / 8 B & data\_offset / point\_count & Offset i liczba próbek, u64.\\
40 / 44 & 4 / 4 B & block\_crc32 / reserved & CRC bloku / zero, u32.\\
\end{longtable}
Dla $N_p$ punktów, $K$ bloków i $D$ bajtów słownika:
\[
 index\_offset=48+D+16N_p,\qquad
 file\_size=48+D+16N_p+48K+4.
\]
Dwa punkty przykładowej metryki: $D=74$, plik 206 B,
dane od 122, indeks od 154, końcowe CRC od 202.

Start kontroluje długości, sumę pliku, klucze, offsety,
kolejność czasów i zgodność min/max; opcjonalnie CRC całego czytanego bloku.
Słownik/indeks pozostają w RAM. Deskryptory są wiązane z bieżącym head
przez LabelSet; próbki czyta się pozycyjnie/binarne wyszukiwanie,
bez wczytania całego Series. Pliki pozostają niezmienne.

\subsection{Manifest i katalog szeregów}
\label{sec:manifest}
MANIFEST ma kopertę: magic \code{RSTMAN04} (8 B), version u16=4,
flags u16=0, json\_len u32; potem JSON UTF-8 (do 64 MiB) i CRC nagłówka
z payloadem. Bez paddingu i dodatkowych bajtów po CRC.
Przykład po pierwszym cut i flush paczki API; crc32=0 jest miejscem
na rzeczywistą sumę segmentu:
\begin{lstlisting}
{
  "format_version": 4, "generation": 1, "wal_checkpoint": 1,
  "segments": [{
    "generation": 1, "file": "segment-00000000000000000001.bin",
    "wal_checkpoint": 1, "size_bytes": 206, "crc32": 0
  }],
  "series": [{
    "metric": "temperature_celsius",
    "labels": [["instance", "1"], ["service", "api"]],
    "sealed_through": 1792843201000
  }]
}
\end{lstlisting}
Checkpoint $N$ oznacza: wszystkie próbki z plików WAL $\leq N$
odwzorowano w segmentach. Recovery pomija te pliki nawet przed ich
usunięciem i czyta pliki $>N$ od offsetu 0.
Nie ma kotwicy, offsetu replay, checkpoint\_lsn ani next\_series\_id.
Generacja head i ID pliku są niezależne.

Katalog/granice pochodzą z freeze, bez późniejszych Series/active.
Szeregi są uporządkowane po kanonicznym kluczu, segmenty po generacji;
nazwy plików są pojedynczymi komponentami ścieżki.
CRC segmentu powtarza jego końcową sumę; JSON nie narzuca kolejności
pól, liczniki dekoduje się bez f64.
Początek: generation=0, wal\_checkpoint=0 i puste tablice, WAL 1.

Walidacja sprawdza wersję, CRC, rozmiary, rosnące generacje/checkpointy,
unikalne klucze i rozłączne zakresy bloków danego klucza.
Każdy Series katalogu ma blok; sealed\_through to maksimum na dysku,
początkowo $H_s=S_s$. Po odbudowie Index przydziela świeże SeriesId.

\subsection{Indeksy w pamięci i plan odczytu}
Index (rys.~\ref{fig:classes}) używa odwróconych, uporządkowanych list ID
dla metryki i par etykiet. Wybór zaczyna od metryki, przecina listy
matcherów i sprawdza pełny klucz. Pusty matcher wybiera brak etykiety;
pustych wartości nie trzeba indeksować.

Deskryptory Series zawierają segment, min/max, offset i liczbę punktów.
Reguła $t>S_s$ zapewnia rozłączne zakresy: deskryptory eliminują zbędne
segmenty, binarne wyszukiwanie wybiera wycinki próbek.
Indeksy odbudowuje silnik, bez zewnętrznej bazy. Ich RAM zależy
od liczby Series, długości etykiet i segmentów; bez kompaktowania
liczba deskryptorów rośnie.

\section{Checkpoint i odzyskiwanie po awarii}
\label{sec:checkpoint}

\subsection{Trwałe opublikowanie generacji}
Jeden koordynator wykonuje poniższy protokół. $g$ numeruje head,
$N$ jest ID zamkniętego pliku WAL; lokalny LSN nie wyznacza checkpointu.

\begin{longtable}{P{27mm}P{129mm}}
\tabhead{Faza & Operacje w kolejności}
Bariera zapisu &
Uzyskać wyłączną barierę; dokończyć rozpoczęte append--head append,
bez czekania na klientów Written. Bez blokad Series wywołać Wal.cut().
Brak nowych append pozwala flusherowi opróżnić bufor, sync starego,
utworzyć nowy i zwrócić $N$. Na tym etapie zapisy czekają na I/O;
odczyty nie pobierają tej bariery.\\
Freeze &
Nadal pod barierą przenieść active do frozen $g$ pod krótkimi
blokadami Series, podnieść $S_s$, zapamiętać katalog po kluczach.
Uruchomić $g+1$; logged\_gen starszej generacji wymusi nowe definicje.
Zwolnić barierę. Następne append trafiają do WAL o numerze $>N$.\\
Segment w tle &
Zapisać segment $g$ z wal\_checkpoint=$N$;
sync pliku $\rightarrow$ rename $\rightarrow$ sync segments/.
Nowe active przyjmują dane, frozen pozostają czytelne.\\
Manifest &
Zapisać/sync MANIFEST.tmp z katalogiem freeze, $g$, $N$ i segmentami;
atomowy rename i sync katalogu danych. To punkt publikacji checkpointu.
Nie dołączać szeregów utworzonych podczas flush.\\
Podmiana head &
Pod krótkimi blokadami Series podmienić frozen na deskryptory.
Nie nadpisywać nowszych active, $H_s$, last\_lsn ani logged\_gen.
ReadView zachowują uchwyty.\\
Usunięcie WAL &
Dopiero teraz Wal.remove\_through($N$); metoda dodatkowo chroni active
i synchronizuje katalog po usunięciu. Zwolnić własne uchwyty RAM.
Żaden plik $>N$ nie podlega temu usunięciu.\\
\end{longtable}

\begin{figure}[htbp]
\centering
\begin{tikzpicture}
\node[anchor=east,font=\small] at (0,0) {Zapisy};
\node[anchor=east,font=\small] at (0,-1.6) {Flush};
\draw[arrow] (.2,0) -- (12,0);
\draw[arrow] (.2,-1.6) -- (12,-1.6);
\draw[fill=black!12,draw=black!40] (2,.45) rectangle (4,-2);
\node[note] at (3,1) {bariera: drain, sync\\cut N + freeze g};
\node[note] at (1,.3) {WAL do N};
\node[note] at (8,.3) {WAL > N, definicje dla g+1};
\draw[very thick,accent] (4.1,-1.6) -- (9.5,-1.6);
\node[note,above] at (6.8,-1.55) {segment g + manifest z N};
\node[note] at (10.8,-2.2) {remove\_through(N)};
\end{tikzpicture}
\caption{Granica checkpointu pokrywa całe pliki WAL.
Pauza obejmuje I/O cut, ale zapis segmentu i manifestu jest poza barierą.}
\label{fig:checkpoint}
\end{figure}

Platforma docelowa: Linux i lokalny system plików z atomowym rename
oraz sync plików/katalogów \cite{linux-sync}. Written nie chroni przed
awarią hosta. Zakończony checkpoint i close opierają się na gwarancjach
fsync/fdatasync, bez ochrony przed fizycznym uszkodzeniem nośnika.

\subsection{Algorytm uruchomienia}
\begin{enumerate}
\item Sprawdzić konfigurację i uzyskać LOCK; ready=false.
W pustym katalogu trwale utworzyć WAL 1 i początkowy manifest z $N=0$.
Przerwaną inicjalizację wolno dokończyć tylko bez danych/segmentów,
z poprawnym pustym WAL. Inny brak manifestu jest błędem, bez resetu.
\item Sprawdzić manifest i wskazane segmenty, odtworzyć Index po
LabelSet, $H_s=S_s$, last\_lsn=0 i nieustawione logged\_gen.
Nie dołączać segmentów spoza manifestu.
\item Użyć poprawki W-01: pomijać WAL $\leq N$ niezależnie od ich obecności.
WAL $N+1$ musi istnieć, choćby pusty; do najwyższego $h$ wymagać ciągłej numeracji.
Czytać od offsetu 0 w kolejności plików/ramek; błędy długości,
CRC, kind i dekodowania obsługuje W-02.
\item Resolver mapuje ID z logu na bieżące ID zwracane przez
ReplaySink.series(LabelSet). Różne ID z tym samym kluczem trafiają
do jednego Series. Próbki wcześniejsze od definicji są przechowywane
w pending i dostarczane po niej.
\item Sink odtwarza próbki chronologicznie w active generacji $g+1$,
sprawdza semantykę i konflikty. Identyczne powtórzenia pomija.
Punkt w zakresie już pokrytym segmentem jest błędem integracji;
nie jest ponownie zapisywany ani cicho gubiony.
Nie ponawia się testów bieżącego zegara ani zmienionego okna $W$.
\item Niepełny koniec tylko ostatniego WAL: truncate do ostatniej
pełnej ramki i sync. Pełna ramka z błędnym CRC/kind/payload zatrzymuje
start także na końcu. Resolver.finish raportuje late/unknown/dropped;
niezerowe dropped lub błąd sinka blokują gotowość (W-06).
\item Jak w Wal::open, utworzyć świeży WAL o ID
$\max(h,N)+1$, z create\_new i sync katalogu; nie dopisywać do starego.
Nowy WalBuffer i watch zaczynają od LSN=0. Index przydziela kolejne
ID ponad największe \emph{bieżące} ID; nie zachowuje ID z logu.
\item Po walidacji usunąć .tmp, sieroty segmentów i WAL $\leq N$.
Uruchomić wątki WAL i gotowość. Bufor replay i pending podlegają
budżetowi pamięci; za mały limit zatrzymuje start.
Odzyskane próbki od razu kwalifikują się do czasowego flush.
\end{enumerate}

\begin{figure}[htbp]
\centering
\begin{tikzpicture}
\node[box,text width=45mm] (lock) at (0,0) {LOCK, manifest\\wal\_checkpoint=N};
\node[box,text width=45mm] (index) at (0,-1.7) {Segmenty po LabelSet\\nowe ID w Index};
\node[box,text width=45mm] (replay) at (0,-3.4)
 {Replay plików > N\\Resolver + ReplaySink};
\node[draw=accent,diamond,aspect=2.8,align=center,font=\small,
 fill=pale,inner sep=2pt] (tail) at (0,-5.2) {Niepełny koniec\\ostatniego WAL?};
\node[box,text width=36mm] (trim) at (6.5,-5.2) {Truncate + sync\\tylko niepełny ogon};
\node[box,text width=45mm] (ready) at (0,-7.1)
 {Nowy WAL h+1, LSN=0\\dropped=0: gotowość};
\node[box,text width=36mm,fill=white] (fail) at (6.5,-1.7)
 {Przerwanie startu\\raport RecoveryReport};
\draw[arrow] (lock) -- (index);
\draw[arrow] (index) -- (replay);
\draw[arrow] (replay) -- (tail);
\draw[arrow] (tail) -- node[note,above]{tak} (trim);
\draw[arrow] (tail) -- node[note,right]{nie} (ready);
\draw[arrow] (trim.south) |- (ready.east);
\draw[arrow] (index) -- node[note,above]{błąd segmentu} (fail);
\draw[arrow] (replay.east) -| node[note,pos=.25,below]
 {CRC / format / orphan} (fail.south);
\end{tikzpicture}
\caption{Recovery stosuje granicę plików i Resolver z brancha.
Nie wymaga trwałego LSN ani utrwalonych head SeriesId.}
\end{figure}

\subsection{Scenariusze awarii i granice gwarancji}
\begin{longtable}{P{44mm}P{112mm}}
\tabhead{Moment awarii & Recovery i znaczenie dla klienta}
Przed append & Brak ramki i sukcesu.\\
Po append, przed head & Pełna zapisana ramka może wrócić, także bez odpowiedzi.\\
None przed write & Paczka może zniknąć; dopuszczona słabsza trwałość.\\
Written przed okresowym sync & Awaria samego procesu zachowuje ramkę;
awaria hosta może ją utracić.\\
Niepełna ostatnia ramka & Tylko niepełny koniec ostatniego pliku
jest odcinany; bez Written dla tej ramki.\\
Podczas cut / przed segmentem & Stary manifest obowiązuje.
Replay obejmuje też rotacje wykonane przez niedokończony cut.\\
Po rename segmentu, przed manifestem & Segment jest sierotą;
replay według starego $N$, bez podwójnego odtwarzania.\\
Po manifeście, przed/podczas usuwania & WAL $\leq N$ jest pomijany,
nawet jeśli jeszcze istnieje; wszystkie $>N$ pozostają do replay.\\
Definicja nigdy nie nadeszła & Resolver raportuje dropped,
a start produkcyjny kończy się błędem, bez cichego sukcesu.\\
\end{longtable}
Błąd segmentu zatrzymuje przyjęcia, zachowując head/WAL i odczyt
do diagnostyki/restartu. I/O WAL: fail-stop.
Po nieopublikowanym freeze $S_s$ może wrócić do poprzedniego checkpointu;
nie narusza to gwarancji Written.

\subsection{Cykl życia procesu}
\begin{figure}[htbp]
\centering
\begin{tikzpicture}
\node[state] (init) at (0,0) {Inicjalizacja\\i recovery};
\node[state] (ready) at (4.8,0) {Gotowy\\readyz=200};
\node[state] (drain) at (9.6,0) {Zamykanie\\readyz=503};
\node[state] (ro) at (4.8,-2) {Odczyt dostępny\\readyz=503};
\node[state,fill=white] (stop) at (9.6,-2) {Proces\\zatrzymany};
\draw[arrow] (init) -- node[note,above]{walidacja OK} (ready);
\draw[arrow] (ready) -- node[note,above]{SIGTERM} (drain);
\draw[arrow] (ready) -- node[note,left]{błąd flush\\danych} (ro);
\draw[arrow] (ro) -- (stop);
\draw[arrow] (drain) -- node[note,right]{close + sync\\lub timeout} (stop);
\draw[arrow] (ready.south east) -- node[note,sloped,above]
  {I/O WAL: fail-stop} (stop.north west);
\end{tikzpicture}
\caption{Diagram stanów procesu. Błąd startu kończy proces bez gotowości;
ze stanu ograniczonego odczytu operator może wykonać restart/zamknięcie.}
\label{fig:lifecycle}
\end{figure}

SIGTERM zamyka przyjęcia i kończy części pamięciowe zadań.
Koordynator kończy rozpoczęty cut przed close i nie rozpoczyna nowego;
cut i close nie mogą ścigać się o zamknięcie flushera.
Po osiągnięciu Written dla oczekujących klientów wywołuje Wal.close()
w zadaniu blokującym: zapisuje także None, wykonuje końcowy sync
i dołącza wątki. Flush danych kończy lub bezpiecznie przerywa,
zachowując pliki replay; pełny flush head nie jest wymagany.
Limit 30 s może skończyć się wymuszonym zatrzymaniem i zwykłymi
gwarancjami None/Written. Drop biblioteki próbuje zamknąć WAL,
ale poprawny shutdown jawnie sprawdza wynik close.

\section{Język i wykonywanie zapytań}
\label{sec:language}

\subsection{Obsługiwana składnia}
Obsługiwany jest nazwany selektor z dokładnymi matcherami, jedna funkcja
czasowa na szereg oraz skalar (np. 1+1 do próby połączenia).
Gramatyka:

\begin{lstlisting}
expr       = selector | function "(" selector "[" duration "]" ")"
           | scalar ;
selector   = metric [ "{" [ matcher { "," matcher } ] "}" ] ;
matcher    = label "=" string ;
function   = "avg_over_time" | "min_over_time" | "max_over_time"
           | "sum_over_time" | "count_over_time" ;
duration   = positive_integer unit ;
unit       = "ms" | "s" | "m" | "h" | "d" | "w" ;
scalar     = number [ "+" number ] ;
\end{lstlisting}

Nazwy: sekcja~\ref{sec:model}; napisy w podwójnych cudzysłowach
z escape JSON; number to skończone f64 o składni JSON.
Whitespace między tokenami jest ignorowany, parser zużywa całe wejście.
Powtórzony matcher i etykieta z prefiksem \code{__} są błędem;
nazwę metryki podaje się wprost. Maksimum tekstu: 16 KiB.

\begin{lstlisting}
temperature_celsius
temperature_celsius{service="api",instance="1"}
avg_over_time(temperature_celsius{service="api"}[5m])
count_over_time(temperature_celsius[30s])
1+1
\end{lstlisting}

Odrzucane: regex, matchery inne niż =, selektor bez nazwy,
samodzielny range vector, podzapytania, offset, @, zagnieżdżenia
i agregacje między Series (np. \code{sum(temperature_celsius)}).
Tokenizator i parser rekurencyjny tworzą AST, bez generatora pełnego
PromQL. Walidacja sprawdza nazwy, dodatnie/mieszczące się okno i typ
wyniku; błąd zawiera offset bajtowy i opis.

\subsection{Czas obliczenia i wybór próbek}
Instant oblicza wyrażenie w $T$, range na siatce
\[
 T_k=start+k\cdot step,\qquad
 k=0,\ldots,\left\lfloor(end-start)/step\right\rfloor .
\]
Wymagania: $step>0$, $start\leq end$, sprawdzana arytmetyka całkowitych
milisekund. Koniec wchodzi do wyniku tylko na siatce.

\begin{longtable}{P{29mm}P{127mm}}
\tabhead{Wyrażenie & Semantyka dla każdego pasującego szeregu}
Selektor &
Najnowsza próbka w $[T-D,T]$, domyślnie $D=5$ min (lookback).
Bez interpolacji, agregacji i markerów staleness; brak próbki pomija Series.
Timestamp wyniku to $T$, nie czas pomiaru; zachowuje \code{__name__}.\\
Funkcja z $R$ &
Okno $(T-R,T]$: lewa granica wykluczona, prawa włączona
\cite{prom-basics}; lookback nie ogranicza okna.
Zachowuje etykiety, usuwa nazwę metryki. Puste okno pomija Series
dla wszystkich funkcji, także count (brak danych nie oznacza zera).\\
\code{avg_over_time} & Średnia arytmetyczna, równa waga próbek.\\
\code{min_over_time}, \code{max_over_time} & Minimum / maksimum.\\
\code{sum_over_time} & Suma w chronologicznym porządku.\\
\code{count_over_time} & Liczba próbek jako wartość liczbowa.\\
\end{longtable}

Funkcje dotyczą próbek liczbowych \cite{prom-functions}.
Obliczenia f64 używają sumowania kompensowanego; średnia = suma/liczność.
Przy przepełnieniu sumy średnią przelicza się ze składników $x_i/n$.
NaN i nieskończoności wyników koduje się jako napisy JSON.

\textbf{Przykład granic:} punkty (czas w s, wartość):
$(0,10),(60,20),(120,30),(180,40)$.
Dla \code{avg_over_time(metric[2m])}:
$T=120\Rightarrow(0,120]\Rightarrow25$;
$T=180\Rightarrow(60,180]\Rightarrow35$;
$T=300\Rightarrow(180,300]\Rightarrow$ brak elementu.
Zakres $[120,190]$, krok 60 s oblicza tylko 120 i 180.
Przeniesienie próbek do segmentu nie zmienia wyniku.

\subsection{Plan wykonania i limity pracy}
\begin{figure}[htbp]
\centering
\begin{tikzpicture}
\node[box] (parse) at (0,0) {Tekst i parametry\\parser + walidacja};
\node[box] (plan) at (5,0) {Plan\\szeregi i zakres};
\node[box] (snap) at (10,0) {Jeden ReadView\\RAM + segmenty};
\node[box] (read) at (10,-2.2) {Odczyt zakresów\\SampleCursor};
\node[box] (eval) at (5,-2.2) {Obliczenia\\dla siatki czasu};
\node[box] (json) at (0,-2.2) {Wektor / macierz\\odpowiedź JSON};
\draw[arrow] (parse) -- (plan);
\draw[arrow] (plan) -- (snap);
\draw[arrow] (snap) -- (read);
\draw[arrow] (read) -- (eval);
\draw[arrow] (eval) -- (json);
\end{tikzpicture}
\caption{Przepływ zapytania. Każdy etap może zakończyć się błędem lub przekroczeniem budżetu.}
\end{figure}

Plan pobiera $[start-D,end]$ dla selektora lub $(start-R,end]$ dla funkcji,
z bezpiecznym przycięciem do dziedziny czasu. ReadView raz kopiuje
potrzebne active, współdzieli frozen; dysk czyta blokowo (I/O 64 KiB).
Przesuwane okno utrzymuje sumę/liczność, a min/max kolejki monotoniczne,
bez ponownego odczytu całego zakresu w każdym kroku.
Nieskończony stan sumy odbudowuje się z bieżącego okna, aby skrajna
wartość po jego opuszczeniu nie zatruwała następnych wyników.

Budżet zapytania: 1000 Series, 2 mln różnych wejściowych próbek,
1 mln punktów wyniku łącznie, 64 MiB i 30 s. RAM obejmuje metadane widoku,
kopie active, okna i wynik przed serializacją; wspólne frozen rozlicza
globalny budżet head, także podtrzymywane tylko przez ReadView.
Odczyt strumieniowy nie znosi limitu próbek.
Przekroczenie oznacza błąd bez częściowego wyniku; timeout sprawdza się
też w pętlach obliczeń/I/O, rozłączenie sygnalizuje anulowanie.

Wyniki są leksykograficzne po kanonicznym kluczu, punkty chronologiczne.
Wartości f64 to napisy, timestampy liczby sekund Unix.
Instant scalar to para czasu/wartości; range scalar to macierz jednego
szeregu o pustych etykietach.

\section{Komunikacja i publiczne API}
\label{sec:http}

Nazwy istniejących tras w \code{src/api.rs} to
\code{POST /api/v1/write_json} (handler \code{write_json})
i \code{GET /api/v1/read} (handler \code{read}). Obecny zapis przyjmuje
tablicę obiektów \code{labels}/\code{samples}, gdzie próbka ma pola
\code{t}/\code{v}; odczyt używa \code{start}, \code{end}, \code{name}
i filtrów etykiet. Poniżej opisano docelowe rozszerzenie kontraktu zapisu
oraz planowane endpointy Prometheusa i diagnostyki.


\subsection{Transport i konwencje}
HTTP/1.1 z keep-alive, UTF-8, JSON \code{application/json};
bez WebSocketów, brokera i skompresowanych żądań.
Profil zapytań przejmuje endpointy, kopertę, wektory i macierze
Prometheusa \cite{prom-api}; nie zapewnia pełnego PromQL/Prometheusa.
Własny zapis nie jest Remote Write ani OpenMetrics.

\subsection{Zapis pojedynczy i paczkowy}
\code{POST /api/v1/write_json} przyjmuje dokładnie jeden z dwóch wariantów:
obiekt pojedynczego pomiaru albo obiekt zawierający wyłącznie niepustą
tablicę \code{points}. Każdy pomiar wymaga pól \code{metric},
\code{timestamp} i \code{value}; brak \code{labels} oznacza pusty zbiór.
\code{timestamp} jest całkowitą liczbą \textbf{milisekund}, a \code{value}
liczbą JSON. W jednej paczce wolno przesłać do 10\,000 punktów.
Maksymalny rozmiar surowego żądania wynosi 4 MiB.
Opcjonalny parametr query \code{durability=none|written}
nadpisuje ustawienie procesu; domyślnie jest to none.
Wartość synced jest nieobsługiwana (400).
Powtórzony parametr, nieznana nazwa albo wartość jest błędem 400.
Tryb dotyczy całej paczki i nie zmienia jej atomowości.

\begin{lstlisting}
POST /api/v1/write_json?durability=written
Content-Type: application/json

{
  "metric": "temperature_celsius",
  "labels": {"service": "api", "instance": "1"},
  "timestamp": 1792843200000,
  "value": 24.5
}
\end{lstlisting}

Wariant paczkowy:
\begin{lstlisting}
{
  "points": [
    {
      "metric": "temperature_celsius",
      "labels": {"service": "api", "instance": "1"},
      "timestamp": 1792843200000,
      "value": 24.5
    },
    {
      "metric": "temperature_celsius",
      "labels": {"service": "api", "instance": "1"},
      "timestamp": 1792843201000,
      "value": 25.5
    }
  ]
}
\end{lstlisting}

Odpowiedź dla pierwszej paczki dwóch nowych punktów po otwarciu,
przy \code{durability=written}:
\begin{lstlisting}
HTTP/1.1 200 OK
Content-Type: application/json

{
  "status": "success",
  "data": {"inserted": 2, "duplicates": 0, "lsn": 143, "durability": "written"}
}
\end{lstlisting}

\code{inserted + duplicates} = liczba punktów wejścia.
LSN to koniec ramki w bieżącym Wal::open, nie trwały identyfikator
żądania; po restarcie może się powtórzyć. Durability potwierdza osiągnięty tryb.
Gwarancje określa sekcja~\ref{sec:wal-runtime}, a reguły duplikatów
i atomowości sekcja~\ref{sec:model}.
Przed enqueue błąd odrzuca całą paczkę; po enqueue timeout/anulowanie/
utrata połączenia nie wycofują zmian. Timeout daje 503/timeout
z niepewnym wynikiem: WAL może osiągnąć wymagany poziom później.

\subsection{Zapytania chwilowe i zakresowe}
Obsługiwane są metody GET oraz POST z
\code{application/x-www-form-urlencoded}. Nie wysyła się JSON jako
ciała zapytania Prometheus. Parametry w danym żądaniu pochodzą albo
z query string, albo z formularza; duplikaty parametrów i mieszanie
obu źródeł są błędem. Całość parametrów nie może przekraczać 32 KiB.

\begin{longtable}{P{36mm}P{40mm}P{80mm}}
\tabhead{Endpoint & Parametry & Znaczenie}
\code{/api/v1/query} & \code{query}; opcjonalnie \code{time},
\code{timeout} & \code{time} domyślnie równe chwili odebrania żądania.\\
\code{/api/v1/query_range} & \code{query}, \code{start}, \code{end},
\code{step}; opcjonalnie \code{timeout} & Obliczenie na siatce czasu;
wszystkie cztery główne parametry są wymagane.\\
\end{longtable}

Czasy przyjmują sekundy Unix z częścią ułamkową lub RFC 3339 ze strefą.
Po konwersji muszą mieścić się w zakresie czasu z sekcji~\ref{sec:model}.
Wymagana jest dokładność co najwyżej milisekundowa; dodatkowe niezerowe
cyfry są odrzucane zamiast cichego zaokrąglania. Krok przyjmuje dodatnią
liczbę sekund, np. \code{0.5}, albo pojedynczą jednostkę jak \code{15s}.
Okno i timeout przyjmują dodatnią liczbę całkowitą z jedną jednostką
z gramatyki \code{duration}. Timeout efektywny jest mniejszą wartością
limitu serwera i parametru klienta.

Nieznane parametry zwracają 400; nie są ignorowane.
Parametry opcjonalne pełnego API Prometheusa, np. \code{lookback_delta}
i \code{limit}, pozostają nieobsługiwane w tym profilu.
Grafana korzysta z konfiguracji nieprzesyłającej tych rozszerzeń.

\begin{lstlisting}
GET /api/v1/query?query=temperature_celsius&time=1792843201

{
  "status": "success",
  "data": {
    "resultType": "vector",
    "result": [
      {
        "metric": {"__name__": "temperature_celsius", "service": "api", "instance": "1"},
        "value": [1792843201, "25.5"]
      }
    ]
  }
}
\end{lstlisting}

Przykład odpowiedzi zakresowej dla tej samej metryki od 1792843200
do 1792843201 z krokiem 1 s:
\begin{lstlisting}
{
  "status": "success",
  "data": {
    "resultType": "matrix",
    "result": [
      {
        "metric": {"__name__": "temperature_celsius", "service": "api", "instance": "1"},
        "values": [[1792843200, "24.5"], [1792843201, "25.5"]]
      }
    ]
  }
}
\end{lstlisting}

Brak danych oznacza \code{result: []} z HTTP 200. Wynik \code{1+1}
w zapytaniu chwilowym ma \code{resultType: "scalar"} oraz
\code{result: [1792843201, "2"]} dla powyższego czasu.
Wartości specjalne wyników obliczeń koduje się jako \code{"NaN"},
\code{"+Inf"} lub \code{"-Inf"}. Wejście zapisu nadal ich nie dopuszcza.

\subsection{Metadane i stan serwera}
\begin{longtable}{P{65mm}P{91mm}}
\tabhead{Metoda i ścieżka & Kontrakt profilu}
GET \code{/api/v1/labels} &
Posortowana lista nazw etykiet oraz \code{__name__} z katalogu.\\
GET \code{/api/v1/label/<name>/values} &
Posortowane unikalne wartości etykiety; dla \code{__name__} nazwy metryk.
Nieistniejąca poprawna nazwa daje pustą listę.\\
GET \code{/api/v1/status/buildinfo} &
Koperta sukcesu z \code{data.version}, \code{data.revision},
\code{data.branch}, \code{data.buildDate} i \code{data.implementation}.
To dane kompilacji TSDB; nie podszywają się pod wersję Prometheusa.\\
GET \code{/healthz} &
200 i \code{ok} jako tekst, jeśli działa pętla obsługi procesu.\\
GET \code{/readyz} &
200 i \code{ready} po recovery i przy sprawnym zapisie;
503 i tekst przy braku gotowości.\\
\end{longtable}

Odpowiedzi metadanych stosują kopertę
\code{status: "success", data: [...]}. Endpointy etykiet nie przyjmują
w tym profilu filtrów czasu ani \code{match[]}; operują na całym katalogu.
Nie obsługuje się \code{/series}, \code{/metadata}, reguł, exemplarów
ani administracyjnego usuwania danych. Dashboard nie wymaga tych operacji.

\subsection{Błędy i przeciążenie}
\label{sec:http-errors}
\begin{longtable}{P{14mm}P{39mm}P{103mm}}
\tabhead{HTTP & errorType & Zastosowanie}
400 & \code{bad_data} & Błędny JSON, parametry, składnia, nazwy,
typy albo nieobsługiwana konstrukcja języka.\\
404 & \code{not_found} & Nieznany endpoint.\\
405 & \code{bad_data} & Niewłaściwa metoda; nagłówek \code{Allow}.\\
409 & \code{conflict} & Różne wartości dla tego samego czasu i szeregu.\\
413 & \code{too_large} & Za duże ciało, tekst zapytania lub liczba punktów paczki.\\
415 & \code{bad_data} & Nieobsługiwany typ treści lub kodowanie wejścia.\\
422 & \code{execution} & Zamknięty zakres, przekroczone okno opóźnienia,
limit szeregów, pamięci lub rozmiaru wykonywanego zapytania.\\
503 & \code{unavailable} / \code{timeout} & Pełna kolejka, zajęte miejsca
zapytań, brak gotowości, timeout albo zatrzymany zapis po błędzie I/O.\\
500 & \code{internal} & Nieoczekiwany błąd lub uszkodzenie podczas odczytu;
szczegóły ścieżek i systemu trafiają do logów.\\
\end{longtable}

\begin{lstlisting}
HTTP/1.1 422 Unprocessable Entity
Content-Type: application/json

{
  "status": "error",
  "errorType": "execution",
  "error": "point belongs to a sealed time range",
  "details": {"reason": "sealed_range", "point_index": 1}
}
\end{lstlisting}

\code{details} jest rozszerzeniem własnego endpointu zapisu; indeks jest
liczony od zera w oryginalnej paczce. Dla zapytań wystarcza standardowa
koperta z \code{status}, \code{errorType} i \code{error}.
Powody odrzucenia out-of-order rozróżnia się jako
\code{sealed_range} i \code{out_of_order_window}.

Globalny semafor dopuszcza do 8 żądań zapisu łącznie: dekodowanych,
oczekujących i przetwarzanych. Adapter rezerwuje miejsce przed
zbuforowaniem ciała. Wewnątrz limitu pula wykonuje domyślnie cztery
operacje równolegle; kolejka przekazuje zadania FIFO, ale nie narzuca
kolejności zakończenia rozłącznych paczek.
Zajęte miejsca skutkują od razu 503 z \code{Retry-After: 1}.
Oczekiwanie na pamięć nie utrzymuje blokad potrzebnych flush.
Limit ośmiu przyjęć obejmuje również żądania oczekujące na Written,
choć nie zajmują już one miejsca w puli części pamięciowej.
Zapytania korzystają z osobnego semafora dwóch miejsc, również bez
nieograniczonego oczekiwania.

Generator ponawia 503 z wykładniczym opóźnieniem od 1 s do 30 s
i losowym odchyleniem. Nie ponawia automatycznie błędów walidacji,
konfliktu i zamkniętego zakresu. Po błędzie połączenia dopuszcza
niepewność wyniku opisaną wcześniej.

\section{Technologie, konfiguracja i środowisko}
\label{sec:technology}

\subsection{Wybór technologii}
\begin{longtable}{P{37mm}P{119mm}}
\tabhead{Wybór & Uzasadnienie, alternatywy i ograniczenia}
Rust 2024 & Kontrola danych i własności pamięci; zamiast C++/Go.
Koszt nauki, konieczne własne limity alokacji.\\
Axum 0.8, Tokio 1 & Routing, ekstraktory i runtime \cite{axum};
alternatywy: Actix Web/Hyper. Blokujące I/O poza wątkami sieci.\\
Serde 1, serde\_json 1 & Typowane DTO zamiast ręcznego JSON;
własne limity i wykrywanie powtórzonych kluczy.\\
time 0.3, Clap 4 & RFC 3339 i jeden sposób konfiguracji CLI;
alternatywy: chrono/plik TOML.\\
std::fs, libc 0.2, własny kodek & Własny silnik jako cel projektu, jawne/testowalne
formaty; zamiast RocksDB/SQLite, kosztem odpowiedzialności za recovery.\\
Arc, RwLock, Mutex, kanały & Blokady Series, bariera head, WAL/watch;
pula pamięciowa i async wait. Bez spinlocków i map lock-free.\\
crc32fast 1, thiserror 2 & Gotowy CRC i typowane błędy zamiast własnego
algorytmu i błędów tekstowych.\\
tracing 0.1, tracing-subscriber 0.3 & Zdarzenia i czasy operacji
zamiast własnego loggera; bez pełnych treści pomiarów.\\
Grafana OSS 13.2.0 & Wykresy, wybór czasu i provisioning
\cite{grafana-provision}; zamiast własnego GUI, zgodność do sprawdzenia.\\
Compose v2, Linux & Sieć, usługi i wolumeny \cite{compose};
zamiast ręcznej instalacji/Kubernetes, bez orkiestracji klastra.\\
cargo test, Tower 0.5, generator Rust/reqwest 0.12 & Testy modułów,
HTTP i obciążenia w jednym stosie, stałe ziarno generatora;
zamiast zewnętrznego frameworka.\\
\end{longtable}

To wybrane linie zgodności, bez zmiany zależności prototypu.
Implementacja utrwala wersje w \code{Cargo.lock}, toolchain
w \code{rust-toolchain.toml}; demo przypina
\code{grafana/grafana:13.2.0} oraz zapisuje digest.
Walidacja: Linux x86-64; inne systemy przez kontenery Linux.

\subsection{Konfiguracja procesu}
Parametry WAL mapują się na WalConfig brancha; limit W-03 i timeout
należą do planowanej integracji. Źródłem konfiguracji są argumenty CLI; brak parametru oznacza wartość
domyślną. Compose przekazuje je jako tablicę \code{command}.
Nie wprowadza się równoległych zmiennych środowiskowych dla tych samych
ustawień. Zmiany wymagają restartu.

\begin{longtable}{P{54mm}P{25mm}P{77mm}}
\tabhead{Parametr & Domyślnie & Znaczenie i walidacja}
\code{--listen} & \code{127.0.0.1:8080} &
Adres HTTP; w kontenerze jawnie \code{0.0.0.0:8080}.\\
\code{--data-dir} & \code{./data} & Katalog dostępny do zapisu,
jeden proces na katalog.\\
\code{--buffer-duration} & \code{30m} &
Czas do zlecenia flush; dodatnia pojedyncza jednostka czasu.\\
\code{--write-workers} & 4 &
Liczba zadań zapisu wykonywanych równolegle, od 1 do 8;
nie zmienia liczby logicznych WAL.\\
\code{--wal-durability} & \code{none} &
None/Written per proces; żądanie może nadpisać.\\
\code{--wal-flush-interval} & \code{100ms} &
Dodatni flush\_interval bez waiterów.\\
\code{--wal-flush-threshold} & 4194304 &
4 MiB, dodatni próg pobudzenia flushera.\\
\code{--wal-buffer-capacity} & 8388608 &
8 MiB początkowej pojemności każdego z dwóch Vec; nie limit.\\
\code{--wal-group-delay-us} & 200 &
commit\_delay w mikrosekundach, nieujemny.\\
\code{--wal-commit-bytes} & 262144 &
256 KiB; od tego rozmiaru pomijane opóźnienie grupy.\\
\code{--wal-sync-interval} & \code{1s} &
Dodatni okres syncera, bez oczekiwania Synced klienta.\\
\code{--wal-writeback} & \code{true} &
Linux sync\_file\_range; false wyłącza optymalizację.\\
\code{--wal-segment-bytes} & 134217728 &
Dodatni miękki próg 128 MiB, sprawdzany przed kolejną porcją.\\
\code{--max-wal-buffer-bytes} & 33554432 &
Planowana W-03: 32 MiB na obie pojemności; co najmniej ich suma
i miejsce na przyjmowaną ramkę, bez nieograniczonej kolejki.\\
\code{--wal-wait-timeout} & \code{30s} &
Dodatni timeout Written w adapterze; bez rollbacku.\\

\code{--max-buffer-bytes} & 536870912 & 512 MiB; dodatni,
co najmniej 65536 B. Wartość 0 wyłącza limit próbek, lecz nie pozostałe limity.\\
\code{--out-of-order-window} & \code{5m} &
Nieujemny, nie większy od czasu buforowania. Nadal obowiązuje $t>S_s$.\\
\code{--max-future-skew} & \code{60s} & Dopuszczalne wyprzedzenie zegara,
nieujemne.\\
\code{--lookback-delta} & \code{5m} & Dodatni zakres wyboru ostatniej próbki.\\
\code{--max-series} & 10000 & Maksymalna liczba szeregów w katalogu;
liczba dodatnia.\\
\code{--max-metadata-bytes} & 67108864 &
64 MiB na katalog i indeksy, z uwzględnieniem kopii katalogu dla checkpointu.\\
\code{--query-timeout} & \code{30s} & Dodatni limit pojedynczego zapytania.\\
\code{--log-level} & \code{info} & Jeden z: error, warn, info, debug, trace.\\
\end{longtable}

Domyślne, stałe limity ochronne z poprzednich sekcji tworzą profil v1:
4 MiB wejścia, 10\,000 punktów paczki, 8 przyjętych żądań zapisu,
dwa zapytania,
64 MiB na zapytanie i 8 MiB całej ramki WAL. Łączny budżet surowych ciał
i zdekodowanych paczek oraz przygotowanych ramek WAL wynosi 64 MiB.
Obejmuje także paczki aktualnie przetwarzane. Pamięć jest rezerwowana
przed dekodowaniem i przy zmianie reprezentacji; przekroczenie budżetu
daje 503 przed enqueue. Osobne 32 MiB WAL obejmuje kopie już przekazane
do bufora i flushera. Odczyt ciała ma timeout 10 s.

Limit metadanych sprawdza się przed rejestracją szeregu oraz przed
utrwaleniem kolejnej paczki, uwzględniając miejsce na indeks następnego
zrzutu. Brak miejsca odrzuca nową paczkę, zamiast przyjmować dane,
których zaplanowany flush przekroczyłby limit. Operator może podnieść
limit po restarcie, co zmienia założenia zużycia RAM.
Wyłączenie limitu bufora nie wyłącza kontroli przepełnień ani błędów alokacji.

\subsection{Sieć, trwałość i granica dostępu}
\label{sec:network}
Topologię, porty, DNS i montowania pokazuje rys.~\ref{fig:deployment}.
Publikacje Compose: \code{127.0.0.1:3000:3000}
i \code{127.0.0.1:8080:8080}.
Baza działa bez root z prawem zapisu do własnego wolumenu.
Grafana wymaga logowania; hasło administratora dostaje przez nieśledzony
plik środowiskowy. API v1 nie ma auth/TLS: granicą jest lokalny host
i zaufana sieć demo; wyjście poza nią wymaga osobnej warstwy dostępu.
Przeglądarka nie łączy się z TSDB, więc CORS nie jest potrzebny.

Compose czeka na \code{/readyz} przed startem generatora.
Po restarcie generator ponawia połączenie, Grafana odzyskuje odczyt
pod stałym adresem. Zatrzymanie kontenerów zachowuje wolumen;
jego usunięcie jest osobną operacją operatora.

\subsection{Diagnostyka}
Logi żądań: request ID, endpoint, wynik, czas, lokalny LSN,
None/Written i liczniki punktów; identyfikator uruchomienia odróżnia
powtarzalne LSN po restartach. Flush: generacja, wal\_checkpoint,
punkty, bajty, czas cut i zapisu segmentu. RecoveryReport:
pliki, rekordy, obcięcie oraz late\_series, unknown\_series, dropped\_samples.
Bez pełnych ciał i etykiet w każdym logu.

Pomiary: RSS, head, rzeczywiste pojemności obu Vec WAL,
backlog\_bytes, waiterzy, write, fdatasync, writeback, pauza cut,
kolejka i flush. Syncer nie udostępnia granicy synced.
Liczniki mogą trafiać okresowo do logów, bez dodatkowego systemu monitoringu.

\section{Wizja interfejsu użytkownika i demonstracja}
\label{sec:ui}

\subsection{Dashboard Grafany}
Dashboard ,,TSDB -- telemetria serwisów'' zawiera wybór jednego serwisu,
selektor czasu, automatyczne odświeżanie co 5 s i cztery panele.
Początkowy zakres to ostatnie 15 min, strefa wyświetlania UTC.
Źródła generatora mają etykiety \code{service} i \code{instance};
różne instancje są osobnymi liniami na wykresie.

\begin{figure}[htbp]
\centering
\begin{tikzpicture}[x=1cm,y=1cm]
\draw[draw=accent,fill=black!2,rounded corners] (0,0) rectangle (14,8);
\node[anchor=west,font=\bfseries] at (0.4,7.55) {TSDB -- telemetria serwisów};
\draw[fill=white,draw=black!30] (0.4,6.6) rectangle (13.6,7.15);
\node[anchor=west,font=\small] at (0.6,6.87)
 {Serwis: api \quad Zakres: ostatnie 15 min \quad UTC \quad Odświeżanie: 5 s};
\foreach \x/\y/\title in {
  0.4/3.5/Temperatura -- pomiary,
  7.2/3.5/Temperatura -- średnia 5 min,
  0.4/0.4/Temperatura -- minimum i maksimum,
  7.2/0.4/Liczba próbek w oknie 1 min} {
  \draw[fill=white,draw=black!30] (\x,\y) rectangle ++(6.4,2.8);
  \node[anchor=west,font=\small] at (\x+0.2,\y+2.5) {\title};
  \draw[black!40] (\x+0.5,\y+0.5) -- ++(5.5,0);
  \draw[black!40] (\x+0.5,\y+0.5) -- ++(0,1.5);
  \draw[accent,thick] plot coordinates {
    (\x+0.6,\y+1.0) (\x+1.2,\y+1.4) (\x+1.8,\y+1.2)
    (\x+2.4,\y+1.7) (\x+3.0,\y+1.4) (\x+3.6,\y+1.8)
    (\x+4.2,\y+1.6) (\x+4.8,\y+1.9) (\x+5.8,\y+1.7)};
  \node[font=\scriptsize] at (\x+3.2,\y+0.2) {czas};
}
\end{tikzpicture}
\caption{Makieta układu dashboardu. Linie są ilustracją, nie wynikami eksperymentu.}
\end{figure}

\begingroup
\small
\begin{longtable}{P{36mm}P{120mm}}
\tabhead{Panel & Wyrażenie projektowanego dashboardu}
Pomiary & \code{temperature_celsius{service="$service"}}\\
Średnia 5 min & \code{avg_over_time(temperature_celsius{service="$service"}[5m])}\\
Minimum i maksimum & Dwa zapytania:
\code{min_over_time(temperature_celsius{service="$service"}[5m])}
oraz \code{max_over_time(temperature_celsius{service="$service"}[5m])}.\\
Liczba próbek & \code{count_over_time(temperature_celsius{service="$service"}[1m])}\\
\end{longtable}
\endgroup

Zmienna \code{service} jest typu \emph{Custom}, z wartościami
\code{api} i \code{worker}. Wyłączone są opcje wielokrotnego wyboru
i ,,All'', aby Grafana nie generowała regexów. Zapytania są zapisane
w trybie \emph{Code}. Minimalny krok panelu to 1 s; Grafana może go
zwiększać dla dłuższego zakresu. Nie używa się makra generującego
funkcje spoza podzbioru.

Brak danych jest przedstawiany jako przerwa lub ,,No data'', bez
łączenia luk wypełnieniem zerami. Błąd API pozostaje widocznym błędem
panelu, a nie pustym poprawnym wynikiem. Panel liczby próbek nie jest
licznikiem wszystkich szeregów ani agregacją między instancjami.

\subsection{Profil standardowego źródła Prometheus}
Provisioning ustawia stały UID, adres bazy i metodę POST.
Poniższy fragment jest projektowaną konfiguracją do użycia podczas
implementacji demo; znaczenie ustawień źródła wynika z dokumentacji
Grafany \cite{grafana-source}.
\begin{lstlisting}
apiVersion: 1
datasources:
  - name: TSDB
    uid: tsdb
    type: prometheus
    access: proxy
    url: http://tsdb:8080
    isDefault: true
    editable: false
    jsonData:
      httpMethod: POST
      timeInterval: 1s
      manageAlerts: false
      disableRecordingRules: true
      prometheusType: Prometheus
      prometheusVersion: 2.0.0
\end{lstlisting}

\code{prometheusVersion: 2.0.0} wybiera konserwatywny profil funkcji
klienta, nie deklaruje wersji silnika TSDB. Prawdziwą wersję zwraca
\code{buildinfo}. Wyłącza się incremental querying, exemplary
i funkcje edytora wymagające nieobsługiwanych endpointów.
Provisioning dashboardu wskazuje katalog plików JSON i UID \code{tsdb}.
Brakującą funkcję Grafany weryfikuje się na przypiętym obrazie,
bez deklarowania zgodności ze wszystkimi wersjami.

\subsection{Scenariusz demonstracji}
\begin{enumerate}
\item Uruchomić Compose z pustym lub istniejącym wolumenem. Poczekać
na gotowość bazy i uruchomienie generatorów.
\item Generatory \code{api} i \code{worker} wysyłają co sekundę temperaturę
dwóch instancji. Jeden profil celowo wprowadza opóźnione punkty.
\item Otworzyć dashboard, zmienić serwis i zakres czasu; porównać
pomiary z agregacjami i liczbą próbek.
\item W profilu demonstracyjnym skrócić buforowanie do 1 min,
a okno opóźnienia do 30 s. Pokazać potwierdzane zapisy i odczyty podczas
flush, niezmienność wcześniejszych danych po podmianie reprezentacji
oraz błąd punktu skierowanego do zamrożonego zakresu.
\item Zrestartować bazę z zachowaniem wolumenu. Pokazać wcześniejsze dane
i dalszy napływ pomiarów.
\end{enumerate}

Profil pokazowy z krótkim buforem nie jest profilem testu zasobów
z lab1. Scenariusz awarii z SIGKILL jest wykonywany w testach
kontrolowanych; sam restart demonstracyjny nie dowodzi gwarancji NFR-04.

\section{Weryfikacja architektury i kryteria odbioru}
\label{sec:verification}

\subsection{Powiązanie z wymaganiami}
\begin{longtable}{P{25mm}P{58mm}P{73mm}}
\tabhead{Wymaganie & Mechanizm & Weryfikacja}
FR-01 & HTTP, walidacja, atomowa ramka WAL. &
T-01: zapis pojedynczy i paczka; błędna paczka nie zmienia danych.\\
FR-02 & Kanoniczny klucz i katalog szeregów. &
T-02: kolejność etykiet, inne wartości, puste etykiety i kolizje hash.\\
FR-03 & Written, bufor WAL i checkpoint. &
T-03, T-10, T-16: replay potwierdzonych Written; osobna ocena None.\\
FR-04 & Siatka czasu i jeden ReadView. &
T-04: lookback, granice, krok niedzielący zakresu i puste wyniki.\\
FR-05 & Parser ograniczonej gramatyki. &
T-05: wszystkie konstrukcje, pełne zużycie wejścia, odrzucane rozszerzenia.\\
FR-06 & Agregacje osobno dla szeregu. &
T-06: pięć funkcji, granice okna, kilka szeregów, tolerancja błędu f64.\\
FR-07 & API Prometheus i provisioning. &
T-07: dashboard i test połączenia na przypiętym obrazie Grafany.\\
FR-08 & Rotacja bloków, budżet pamięci i flush w tle. &
T-08, T-14--15: timer, nowe zapisy podczas flush i pamięć utrzymywana przez odczyt.\\
FR-09 & $W$, $H_s$, $S_s$. &
T-09: granice równości, flush, zmiana konfiguracji i restart.\\
NFR-01--03 & Paczki, indeksy, limity i kontrola alokacji. &
P-01--03: testy HTTP, przepustowości i RSS z opisanym profilem.\\
NFR-04 & Recovery i kolejność publikacji plików. &
T-10 i P-04: wymuszone awarie i czas gotowości.\\
NFR-05 & Typowane błędy i walidacja. &
T-11: macierz błędnych żądań bez zakończenia procesu.\\
DA-02--04, DA-08 & Index, blokady Series, bariera rotacji i jeden WAL. &
T-12--15 i P-05: równoległość, brak zakleszczeń i poprawność widoków.\\
DA-07, DA-09 & Flusher, syncer, granice prefiksów i tryby per żądanie. &
T-16--20 i P-06: przyczynowość definicji, rotacja, checkpoint i błędy I/O.\\
\end{longtable}

\subsection{Scenariusze poprawności}
\begin{description}[style=multiline,leftmargin=16mm,font=\bfseries]
\item[T-01:] Jeden punkt, wieloszeregowa paczka i 10\,000 punktów:
odczyt po sukcesie widzi całość. Błędny punkt nie zmienia próbek/katalogu;
niewidoczne luki ID dozwolone.
\item[T-02:] Kolejność etykiet nie zmienia Series, wartość zmienia;
puste usuwane, powtórzone klucze JSON/nazwy zastrzeżone odrzucane,
kolizja hash nie scala kluczy. Adapter dodaje wewnętrzne \_\_name\_\_.
\item[T-03:] Restart z WAL, samymi segmentami/pustym nowym WAL i stanem
mieszanym: te same próbki/etykiety/$S_s$, mimo zmiany head ID.
Nowy plik ma rosnące ID, pierwszy LSN jest liczony od 0.
\item[T-04:] Instant/range: start=end, krok=0, odwrócony zakres,
koniec poza siatką, dokładna granica lookback, wiele Series/brak danych.
\item[T-05:] Parser: funkcje, escape, whitespace, 1+1;
błędna jednostka, trailing input, regex i nieobsługiwana funkcja.
\item[T-06:] Ręczne wyniki, puste/brzegowe $(T-R,T]$, skrajne wejścia
i nieskończony wynik; RAM/dysk/mieszany widok z tolerancją
$10^{-9}\max(1,|wynik_{\mathrm{ref}}|)$.
\item[T-07:] Przypięta Grafana: provisioning, wszystkie panele,
wybór czasu/serwisu, przerwy, restart i test połączenia.
\item[T-08:] Flush czasowy/pamięciowy przy odczycie: punkt raz,
bez pełnej kopii RAM; za mały budżet paczki nie zapętla flush.
\item[T-09:] Równość i obie strony $H-W$, $S$, duplikaty/konflikty,
ponowna walidacja po cut/freeze. Zmiana $W$ nie odrzuca replay.
\item[T-10:] $\geq20$ awarii procesu w Written, z osobnym dziennikiem
potwierdzeń klienta: środek ramki, po write/sync, podczas cut,
rename segmentu/manifestu, przed/po usuwaniu WAL.
Potwierdzone Written wracają; None i paczki bez odpowiedzi mogą
wrócić w całości lub być nieobecne. Drop/close nie zastępuje SIGKILL.
\item[T-11:] Błędne typy/NaN, rozmiar, parametry (w tym synced),
pełne przyjęcia/zapytania i timeout Written: właściwy kod/opis,
dalsza diagnostyka. Timeout nie wycofuje publikacji.
\item[T-12:] Co najmniej dwa zapisy rozłącznych Series.
Wstrzymać flusher po head append klienta Written: inny None,
także do tego Series, oraz odczyt mogą skończyć.
Konkurencyjne duplikaty/konflikty wspólnego Series bez utraty próbek.
\item[T-13:] Paczki tych samych Series w odwrotnej kolejności wejścia
bez deadlock; odczyt widzi całość albo stan sprzed paczki.
Rejestracja jednego klucza daje jeden Series; zmiana wersji katalogu
podczas snapshot wymusza ponowienie.
\item[T-14:] Wstrzymać flush po cut/freeze: nowy active działa,
obie generacje czytelne. Pierwszy zapis istniejącego Series w nowym $g$
ponownie zawiera definicję; usunięcie WAL $\leq N$ nie psuje restartu.
Powtórzyć dla Series utworzonego podczas flush.
\item[T-15:] Wolny flush, active bliski limitu i ReadView starego bloku:
RAM rozliczony do ostatniego uchwytu, 503 przed append,
bez kolejki frozen i blokowania zwalniania pamięci.
\item[T-16:] Przy zatrzymanym flusherze None odpowiada,
Written czeka mimo widocznego head; write zwalnia Written.
Okresowy sync nie tworzy dodatkowego trybu/warunku odpowiedzi.
Duplikat None ponowiony jako Written rzeczywiście czeka.
\item[T-17:] Powtórne definicje i różne ID tego samego LabelSet mapują się
na jeden Series. Próbki przed definicją wracają po jej dostarczeniu.
Resolver raportuje definicję, która nigdy nie nadeszła, a gotowość
produkcyjna jest wtedy zablokowana.
\item[T-18:] Written nie zależy od długiego flush\_interval:
waiter budzi flusher. Postęp watch przed subskrypcją, krótki write,
anulowanie i zamknięty kanał nie powodują utraty pobudzenia/fałszywego sukcesu.
Rotacja z równoległym sync zachowuje ważny Arc<File>.
\item[T-19:] Wstrzymany head append opóźnia barierę cut, waiter Written nie.
Cut nie obejmuje nowych append po zwolnieniu bariery.
Po manifeście pozostawione WAL $\leq N$ są pomijane.
Rotacje rozmiarowe wewnątrz generacji i przekroczenie miękkiego progu
o porcję nie zmieniają replay; brak dzielenia ramki.
\item[T-20:] Budżet active+spare i scratch wymusza backpressure przed append,
bez części paczki. Błąd write/sync/writeback kończy proces bez sukcesu
nieosiągniętych Written. Jawny close zapisuje również None, synchronizuje
i odrzuca późniejsze append. Wcześniejsze Written wracają przy sprawnym nośniku.
\end{description}

Kodeki: zgodność bajtowa RecordEncoder, 9-bajtowy nagłówek,
len payloadu, CRC(kind || len || payload), każde ucięcie małej ramki,
nadmierne liczniki/offsety i pełna błędna ramka na końcu (błąd startu).
Bariery wymuszają przeploty. Sprawdzić puste pliki, luki po $N$,
replay bez usunięcia starego WAL, restart LSN=0 i brak trwałych head ID.

\subsection{Profil pomiarów wydajności i zasobów}
Raport podaje CPU/rdzenie, RAM, dysk/system plików, OS, toolchain,
commit, konfigurację i współbieżność klienta. Referencja:
Linux x86-64, $\geq4$ rdzenie, 8 GiB RAM, lokalny SSD bez innego
obciążenia, build release; parametry nie zastępują pomiaru.

\begin{longtable}{P{20mm}P{73mm}P{63mm}}
\tabhead{Test & Przebieg & Cel z lab1}
P-01 & Po 30 s rozgrzewki co najmniej 5 min pomiaru. Osobne scenariusze
zapisu jednego punktu, paczki 1000 punktów, odczytu 1000 wyników i agregacji.
Profil bazowy Written; czas HTTP obejmuje kolejkę, publikację head
i oczekiwanie na write. &
p95 zapisu punktu $\leq100$ ms; paczki $\leq500$ ms;
odczytu $\leq200$ ms; agregacji z RAM $\leq500$ ms.\\
P-02 & 1000 szeregów o stałych etykietach, paczki po 1000 punktów,
4 współbieżnych klientów, Written. Co najmniej godzina, by objąć flush;
raportować także najgorsze okna 1 s i 10 s. &
Średnio $\geq10\,000$ nowych potwierdzonych punktów/s,
z uwzględnieniem rotacji, działania flush i błędów.\\
P-03 & 10\,000 punktów/s, 1000 szeregów, buforowanie 30 min,
domyślne limity, minimum 65 min. Bazowy scenariusz bez zapytań,
dodatkowo wariant z dwoma zapytaniami. &
RSS procesu $\leq1$ GiB, w tym szczyt przy flush i realokacji.
Raportować rzeczywisty czas przebywania punktów w RAM.\\
P-04 & Milion potwierdzonych punktów; osobno stan z przewagą WAL
i stan po flush. Pomiar od startu procesu do pierwszego 200 z readyz;
osobno cache zimny i ciepły. &
Gotowość w $\leq10$ s i brak utraty danych w T-10.\\
P-05 & Warianty 1, 2, 4 i 8 zadań zapisu i klientów.
Osobno rozłączne szeregi oraz wszystkie paczki kierowane do jednego Series.
Pomiar obejmuje co najmniej jeden flush. &
Ocena rzeczywistej równoległości, kosztu blokad i wspólnego WAL;
bez założenia liniowego przyspieszenia.\\
P-06 & Osobne przebiegi None i Written dla tych samych paczek.
Zmiana opóźnienia grupy i obserwacja opóźnień flushera/syncera. &
Raportować p95, przepustowość i bajty oczekujące oddzielnie;
None nie zalicza kryterium trwałej przepustowości.\\
\end{longtable}

P-01: selektor jednego Series, 1000 chwil co 1 s i próbki co 1 s;
agregacja tego Series, okno 5 min, ta sama siatka, potrzebne dane w RAM.
P-02/03: 1000 Series po 10 punktów/s, najwyżej cztery krótkie etykiety;
bazowo 0\% out-of-order, dodatkowo 1\%/10\% z opóźnieniem do 5 min.
Odrzuceń po zamknięciu nie wlicza się do przepustowości.

Surowe próbki z 30 min: $10\,000\cdot1800=18$ mln,
$18\,000\,000\cdot16=288\,000\,000$ B $\approx274{,}7$ MiB.
Dochodzi zapas pojemności, active/frozen, uchwyty, metadane, przyjęcia,
scratch, HTTP i zapytania. Limit head 512 MiB daje margines,
\textbf{nie dowodzi RSS $\leq1$ GiB}. Mierzy się też wspólny WAL,
oczekiwanie na Series, barierę z I/O cut i osobno dyskowy flush segmentu poza barierą.
Nieosiągnięty cel raportuje się z wynikiem i wąskim gardłem.

\subsection{Ryzyka i warunki akceptacji projektu}
\begin{longtable}{P{37mm}P{61mm}P{58mm}}
\tabhead{Ryzyko & Sposób ograniczenia & Obserwowany wskaźnik}
Flush zużywa zapas pamięci & Rotacja przy 80\% budżetu, jeden flush
w tle, rozliczanie bloków podtrzymywanych przez odczyty i backpressure. &
Czas flush i bariery, pamięć obu generacji, 503 i p95 zapisu.\\
Konkurencja o blokady i WAL & Krótkie blokady Series, enqueue bez I/O
i oczekiwanie na Written dopiero po publikacji. &
T-12--13, P-05--06, czas części pamięciowej, wait Written i pauzy cut.\\
Koszt out-of-order & Wektor uporządkowany, jawny limit $W$ i zamknięta
granica; test udziału opóźnionych danych. &
Czas wstawienia i liczba przesuniętych próbek.\\
Wzrost indeksów i plików & Limit katalogu i metadanych; jawny brak
retencji i kompaktowania w v1. &
Liczba szeregów, segmentów, bajty metadanych i dysku.\\
Niepełna zgodność Grafany & Ustalony obraz, tryb Code, ograniczone
zmienne i wczesny test rzeczywistych żądań. &
T-07 oraz lista wywołanych endpointów i parametrów.\\
Błąd checkpointu & Manifest jako jedyne źródło prawdy, sync przed rename,
WAL usuwany na końcu. &
T-10 i testy uszkodzeń plików.\\
\end{longtable}

Akceptacja projektu architektury wymaga spójności kontraktów i formatów,
pokrycia wszystkich obszarów wytycznych lab2 oraz prześledzenia scenariuszy
awarii. Akceptacja późniejszej implementacji wymaga wykonania testów
i pomiarów; niniejszy dokument nie oznacza ich zaliczenia.

\appendix
\section{Słownik}
\begin{longtable}{P{32mm}P{124mm}}
\tabhead{Pojęcie & Znaczenie}
Head / Index / Series & Pamięciowy katalog i active/frozen /
jeden katalog kluczy i uchwytów / szereg z własną blokadą i blokami.\\
LabelSet / SeriesId & Kanoniczne etykiety z \_\_name\_\_ /
lokalne ID head/WAL, odtwarzane przez mapowanie po kluczu.\\
WAL / LSN & Bufor logu przed head append /
liczba bajtów przyjętych w bieżącym Wal::open, reset po restarcie.\\
Flusher / Syncer & Zapisuje porcje, publikuje Written i obsługuje cut /
okresowo sync\_data, bez potwierdzenia per żądanie.\\
Cut / Checkpoint & Drain, sync i rotacja WAL zwracające zamknięty plik $N$ /
trwały manifest z danymi wszystkich WAL $\leq N$.\\
Generacja / logged\_gen & Grupa wspólnie zamrażanych bloków /
znacznik, czy Series ma definicję w aktualnej generacji WAL.\\
Segment / Manifest & Niezmienny plik kluczy, próbek i indeksu /
spis segmentów, klucze, granice zamknięcia i wal\_checkpoint.\\
Freeze / Flush & Przeniesienie active do frozen po cut /
zapis segmentu i manifestu w tle, potem remove\_through.\\
Resolver / ReplaySink & Mapowanie ID i pending late series /
odbudowa Index/head po LabelSet.\\
Out-of-order / $S_s$ & Próbka starsza od najnowszej /
granica zamknięcia: nowy czas musi być większy.\\
Lookback / Provisioning / RSS & Maksymalny wiek próbki selektora /
deklaratywna konfiguracja Grafany / rezydentna pamięć procesu.\\
\end{longtable}

\section{Budowanie dokumentu}
Źródło jest samodzielnym plikiem LaTeX. Diagramy i makieta są zapisane
w TikZ, bez zależności od zewnętrznych obrazów lub usług.
Wymagane są pdfLaTeX, latexmk oraz pakiety użyte w preambule,
w szczególności obsługa języka polskiego i TikZ.
Z katalogu głównego repozytorium:
\begin{lstlisting}
latexmk -pdf -interaction=nonstopmode -halt-on-error \
  -outdir=target/lab2-doc \
  pz/lab2/lab2-Blonowski-Zaborski.tex
\end{lstlisting}
Pliki pomocnicze trafiają do ignorowanego katalogu \code{target/lab2-doc}.
PDF do oddania jest kopią wyniku kompilacji umieszczoną obok źródła.
Po zmianach należy sprawdzić ostrzeżenia składu, odwołania, spis treści
oraz czytelność tabel i diagramów.

\begingroup
\small
\begin{thebibliography}{16}
\raggedright
\setlength{\itemsep}{0pt}
\addcontentsline{toc}{section}{Źródła}
\bibitem{lab1}
H. Błonowski, J. Zaborski,
\emph{Etap 1: Specyfikacja wymagań i zakres pracy}, wersja 1.1,
21.10.2026, dokument projektu w \code{pz/lab1}.
\bibitem{wal23}
Repozytorium projektu, gałąź \code{23-wal-improvements},
commit \code{e1354b9}, pliki \code{src/storage/wal},
\code{src/storage/series.rs}, \code{src/db.rs}; odczyt lokalny 02.11.2026.
\url{https://github.com/janzaborski/rust-tsdb/tree/e1354b9/src/storage/wal}.
\bibitem{gorilla}
T. Pelkonen, S. Franklin, P. Cavallaro, Q. Huang, J. Meza,
J. Teller, K. Veeraraghavan,
\emph{Gorilla: A Fast, Scalable, In-Memory Time Series Database},
Proceedings of the VLDB Endowment, 8(12), s. 1816--1827, 2015,
sekcje 4.2--4.3,
\url{https://www.vldb.org/pvldb/vol8/p1816-teller.pdf}.
\bibitem{prom-api}
Prometheus, \emph{HTTP API},
\url{https://prometheus.io/docs/prometheus/latest/querying/api/}.
\bibitem{prom-basics}
Prometheus, \emph{Querying basics},
\url{https://prometheus.io/docs/prometheus/latest/querying/basics/}.
\bibitem{prom-functions}
Prometheus, \emph{Query functions},
\url{https://prometheus.io/docs/prometheus/latest/querying/functions/}.
\bibitem{grafana-source}
Grafana Labs, \emph{Configure the Prometheus data source},
\url{https://grafana.com/docs/grafana/latest/datasources/prometheus/configure/}.
\bibitem{grafana-provision}
Grafana Labs, \emph{Provision Grafana},
\url{https://grafana.com/docs/grafana/latest/administration/provisioning/}.
\bibitem{axum}
Projekt Axum, \emph{Dokumentacja biblioteki},
\url{https://docs.rs/axum/latest/axum/}.
\bibitem{rust-file}
Projekt Rust, \emph{std::fs::File},
\url{https://doc.rust-lang.org/std/fs/struct.File.html}.
\bibitem{crc}
Projekt crc32fast, \emph{Dokumentacja biblioteki},
\url{https://docs.rs/crc32fast/latest/crc32fast/}.
\bibitem{compose}
Docker, \emph{How Compose works},
\url{https://docs.docker.com/compose/intro/compose-application-model/}.
\bibitem{linux-write}
Linux man-pages, \emph{write(2)},
\url{https://man7.org/linux/man-pages/man2/write.2.html}.
\bibitem{linux-sync}
Linux man-pages, \emph{fsync(2), fdatasync(2)},
\url{https://man7.org/linux/man-pages/man2/fsync.2.html}.
\bibitem{linux-range}
Linux man-pages, \emph{sync\_file\_range(2)},
\url{https://man7.org/linux/man-pages/man2/sync_file_range.2.html}.
\bibitem{tokio-watch}
Projekt Tokio, \emph{tokio::sync::watch},
\url{https://docs.rs/tokio/latest/tokio/sync/watch/index.html}.
\end{thebibliography}
\noindent Dokumentację zewnętrzną sprawdzono 24.10.2026.
Zachowanie własnego silnika określa niniejsza specyfikacja; odnośniki
\emph{latest} służą jako kontekst, a zgodność integracji sprawdza się
z przypiętymi wersjami używanymi w demonstracji.
\endgroup

\end{document}
